\documentclass[12pt,UTF-8,openany,fontset=none]{ctexbook}
% 导入公用配置
% 诗歌类导言区
\usepackage{titlesec}
\usepackage{xeCJK}
\usepackage{fontspec,xunicode,xltxtra}
\usepackage{xpinyin}
\usepackage{tcolorbox}
\usepackage{geometry}
\usepackage{indentfirst}
\usepackage{pifont}
\usepackage[perpage,symbol*]{footmisc}
\usepackage{eso-pic}
\usepackage{graphicx} 
\usepackage{tikz}

\ExplSyntaxOn
\cs_set_protected:Npn \__xpinyin_make_pinyin_box:nnn #1#2#3
  {
    \__xpinyin_leavevmode:
    \hbox_overlap_right:n
      {
        \hbox_set:Nn \l__xpinyin_tmpa_box
          { \__xpinyin_CJKsymbol_hook: \__xpinyin_save_CJKsymbol:n {#2} }
        \hbox_set:Nn \l__xpinyin_tmpb_box
          {
            \l__xpinyin_pinyin_box_hook_tl
            \__xpinyin_select_font:
            \l__xpinyin_format_tl
            \clist_if_exist:cT { c__xpinyin_multiple_ #1 _clist }
              { \l__xpinyin_multiple_tl }
            {#3}
          }
        % 原第 164–171 行的 \dim_compare 和 \box_resize_to_wd_and_ht 已删除
        \box_move_up:nn { \l__xpinyin_vsep_tl }
          {
            \hbox_to_wd:nn { \box_wd:N \l__xpinyin_tmpa_box }
              { \tex_hss:D \box_use_drop:N \l__xpinyin_tmpb_box \tex_hss:D }
          }
      }
  }
\ExplSyntaxOff

% 设置背景图
\newcommand{\SetPageBackground}[1]{%
  \AddToShipoutPictureBG*{%
    \AtPageLowerLeft{%
    \begin{tikzpicture}[opacity=0.75]
        \node[inner sep=0pt, anchor=south west] at (0,0)
        {\includegraphics[width=\paperwidth,height=\paperheight,
                            keepaspectratio]{#1}};
    \end{tikzpicture}%
    }%
  }%
}

% 定义拼音使用的缺省字体
\setmainfont{Mona Sans Light}
% 定义 Century Schoolbook 字体
\newfontfamily\schoolbookfont{Century Schoolbook}
% 定义中文字体
\setCJKmainfont[BoldFont=STZhongsong]{KaiTi}
% 定义隶书字体族
\setCJKfamilyfont{beishu}{YouYuan} % 定义背诵框字体
\newcommand{\beishu}{\CJKfamily{beishu}} % 创建一个 \beishu 命令，方便调用
\xeCJKDeclareCharClass{CJK}{`0 -> `9}
\xeCJKsetup{AllowBreakBetweenPuncts=true}
\DefineFNsymbols{circled}{{\ding{192}}{\ding{193}}{\ding{194}}{\ding{195}}{\ding{196}}{\ding{197}}{\ding{198}}{\ding{199}}{\ding{200}}{\ding{201}}}
\setfnsymbol{circled}
\xpinyinsetup{ratio=0.48,hsep={.5em plus .1em minus .1em},vsep={1em}}
\setpinyin{一}{yi4}
\setpinyin{尽}{jin4}
\setpinyin{地}{di4}
\setpinyin{长}{chang2}
\setpinyin{朝}{zhao1}
\setpinyin{罗}{luo2}
\setpinyin{铺}{pu1}
\setpinyin{泊}{bo2}
\setpinyin{似}{si4}


% 定义背诵框环境
\newtcolorbox{beisongbox}{
    colback=yellow!10!white,   % 背景色
    colframe=orange!60!black,  % 边框色
    arc=2pt,                   % 圆角半径
    boxrule=0.6pt,             % 边框线粗细
    halign=center,             % 内容水平居中
    valign=center,             % 内容垂直居中
    nobeforeafter,             % 去掉框前后的额外间距
    boxsep=0pt,                % 内边距
    left=2pt,                  % 左内边距
    right=2pt,                 % 右内边距
    top=2pt,                   % 上内边距
    bottom=2pt,                % 下内边距
    fontupper=\beishu\zihao{3}\color{red!70!black}\ziju{-0.0}, % 框内字体、字号等
}
% 背诵框命令
\newcommand{\beisong}[1]{
    \begin{textblock*}{30mm}(15mm, 13mm)   % 宽 30mm，左上角在 (15mm, 13mm)  
    \begin{beisongbox}#1\end{beisongbox}
\end{textblock*}
}

% 章节标题字体格式
\titleformat{\chapter}{\centering\zihao{0}\bfseries}{}{0pt}{}
\titleformat{\section}{\zihao{-2}\bfseries}{ }{0pt}{}
\author{}
\date{}


\geometry{a4paper,left=1.4cm,right=1.4cm,top=2.3cm,bottom=2.3cm}
\renewcommand{\footnotesize}{\fontsize{10.5pt}{12pt}\selectfont}
% 定义中文字体
\setCJKmainfont[BoldFont=STZhongsong]{方正书宋简体}  % Source Han Serif SC
\xeCJKsetup{AutoFallBack=true}
\setCJKfallbackfamilyfont{\CJKrmdefault}{Source Han Serif SC}

\renewcommand{\thechapter}{}                    % 章不显示编号
\renewcommand{\thesection}{\arabic{chapter}.\arabic{section}}  % 节号仍为 1.1 / 2.1
\title{\zihao{0} \bfseries 中学古诗词集萃}
\setlength{\lineskip}{24pt}
\setlength{\parskip}{6pt}
\author{}
\date{}
\makeatletter
\renewcommand{\CTEXthechapter}{}
\makeatother
\makeatletter
\RenewDocumentCommand{\section}{s o m}{%
  \IfBooleanTF{#1}{%
    \addcontentsline{toc}{section}{#3}%
  }{%
    \refstepcounter{section}%
    \addcontentsline{toc}{section}{\protect\numberline{\thesection}#3}%
  }%
}
\makeatother

% 正文
\begin{document}
\maketitle
\tableofcontents
\newpage

\chapter{第一级}
\clearpage

\section{卜算子·咏梅}

\beisong{泛读}

\begin{center}
    {\zihao{-0}\bfseries 卜算子·咏梅}

    \vspace{18pt}

    \begin{large}\xuwen
        〔毛泽东〕
    \end{large}

\end{center}

    \vspace{18pt}  % 起首

\begin{center}
    \begin{Large}

        风雨送春归，

        飞雪迎春到。

        已是悬崖百丈冰，

        犹有花枝俏。

        \vspace{18pt}  % 分阙

        俏也不争春，

        只把春来报。

        待到山花烂漫时

        她在丛中笑。

    \end{Large}

\end{center}

\vspace{8pt}
\clearpage

\section{菩萨蛮·大柏地}

\beisong{泛读}

\begin{center}
    {\zihao{-0}\bfseries 菩萨蛮·大柏地}

    \vspace{18pt}

    \begin{large}\xuwen
        〔毛泽东〕
    \end{large}

\end{center}

    \vspace{18pt}  % 起首

\begin{center}
    \begin{Large}

        赤橙黄绿青蓝紫，

        谁持彩练当空舞？

        雨后复斜阳，

        关山阵阵苍。

        \vspace{18pt}  % 分阙

        当年鏖战急，

        弹洞前村壁。

        装点此关山，

        今朝更好看。

    \end{Large}

\end{center}

\vspace{8pt}
\clearpage

\section{三五七言}

\beisong{泛读}

\begin{center}
    {\zihao{-0}\bfseries 三五七言}

    \vspace{18pt}

    \begin{large}\xuwen
        〔李白〕
    \end{large}

\end{center}

    \vspace{18pt}  % 起首

\begin{center}
    \begin{Large}

        秋风清，秋月明，

        落叶聚还散，寒鸦栖复惊。

        相思相见知何日？此时此夜难为情！

    \end{Large}

\end{center}

\vspace{8pt}
\clearpage

\section{闻王昌龄左迁龙标遥有此寄}

\beisong{背诵}

\begin{center}
    {\zihao{-0}\bfseries 闻王昌龄左迁龙标遥有此寄}

    \vspace{18pt}

    \begin{large}\xuwen
        〔李白〕
    \end{large}

\end{center}

    \vspace{18pt}  % 起首

\begin{center}
    \begin{Large}

        杨花落尽子规啼，闻道龙标过五溪。

        我寄愁心与明月，随风直到夜郎西。

    \end{Large}

\end{center}

\vspace{8pt}
\clearpage

\section{次北固山下}

\beisong{背诵}

\begin{center}
    {\zihao{-0}\bfseries 次北固山下}

    \vspace{18pt}

    \begin{large}\xuwen
        〔王湾〕
    \end{large}

\end{center}

    \vspace{18pt}  % 起首

\begin{center}
    \begin{Large}

        客路青山外，行舟绿水前。

        潮平两岸阔，风正一帆悬。

        海日生残夜，江春入旧年。

        乡书何处达，归雁洛阳边。

    \end{Large}

\end{center}

\vspace{8pt}
\clearpage

\section{天净沙·秋思}

\beisong{背诵}

\begin{center}
    {\zihao{-0}\bfseries 天净沙·秋思}

    \vspace{18pt}

    \begin{large}\xuwen
        〔马致远〕
    \end{large}

\end{center}

    \vspace{18pt}  % 起首

\begin{center}
    \begin{Large}

        枯藤老树昏鸦，

        小桥流水人家，

        古道西风瘦马。

        夕阳西下，\hspace{25pt}

        断肠人在天涯。

    \end{Large}

\end{center}

\vspace{8pt}
\clearpage

\section{清平乐·村居}

\beisong{背诵}

\begin{center}
    {\zihao{-0}\bfseries 清平乐·村居}

    \vspace{18pt}

    \begin{large}\xuwen
        〔辛弃疾〕
    \end{large}

\end{center}

    \vspace{18pt}  % 起首

\begin{center}
    \begin{Large}

        茅檐低小，溪上青青草。

        醉里吴音相媚好，白发谁家翁媪？

        \vspace{18pt}  % 分阙

        大儿锄豆溪东，中儿正织鸡笼。

        最喜小儿亡赖，溪头卧剥莲蓬。

    \end{Large}

\end{center}

\vspace{8pt}
\clearpage

\section{西江月·夜行黄沙道中}

\beisong{背诵}

\begin{center}
    {\zihao{-0}\bfseries 西江月·夜行黄沙道中}

    \vspace{18pt}

    \begin{large}\xuwen
        〔辛弃疾〕
    \end{large}

\end{center}

    \vspace{18pt}  % 起首

\begin{center}
    \begin{Large}

        明月别枝惊鹊，清风半夜鸣蝉。

        稻花香里说丰年，听取蛙声一片。

        \vspace{18pt}  % 分阙

        七八个星天外，两三点雨山前。

        旧时茅店社林边，路转溪桥忽见。

    \end{Large}

\end{center}

\vspace{8pt}
\clearpage

\section{相思}

\beisong{泛读}

\begin{center}
    {\zihao{-0}\bfseries 相思}

    \vspace{18pt}

    \begin{large}\xuwen
        〔王维〕
    \end{large}

\end{center}

    \vspace{18pt}

\begin{center}
    \begin{Large}

        红豆生南国，春来发几枝？

        愿君多采撷，此物最相思。

    \end{Large}

\end{center}

\vspace{8pt}
\clearpage

\section{秋夕}

\beisong{泛读}

\begin{center}
    {\zihao{-0}\bfseries 秋夕}

    \vspace{18pt}

    \begin{large}\xuwen
        〔杜牧〕
    \end{large}

\end{center}

    \vspace{18pt}

\begin{center}
    \begin{Large}

        银烛秋光冷画屏，轻罗小扇扑流萤。

        天阶夜色凉如水，坐看牵牛织女星。

    \end{Large}

\end{center}

\vspace{8pt}
\clearpage

\section{月夜}

\beisong{背诵}

\begin{center}
    {\zihao{-0}\bfseries 月夜}

    \vspace{18pt}

    \begin{large}\xuwen
        〔刘方平〕
    \end{large}

\end{center}

    \vspace{18pt}

\begin{center}
    \begin{Large}

        更深月色半人家，北斗阑干南斗斜。

        今夜偏知春气暖，虫声新透绿窗纱。

    \end{Large}

\end{center}

\vspace{8pt}
\clearpage

\chapter{第二级}
\clearpage

\section{水调歌头·重上井冈山}

\beisong{精读}

\begin{center}
    {\zihao{-0}\bfseries 水调歌头·重上井冈山}

    \vspace{18pt}

    \begin{large}\xuwen
        〔毛泽东〕
    \end{large}

\end{center}

    \vspace{18pt}  % 起首

\begin{center}
    \begin{Large}

        久有凌云志，

        重上井冈山。

        千里来寻故地，

        旧貌变新颜。

        到处莺歌燕舞，

        更有潺潺流水，

        高路入云端。

        过了黄洋界，

        险处不须看。

        \vspace{18pt}  % 分阙

        风雷动，

        旌旗奋，

        是人寰。

        三十八年过去，

        弹指一挥间。

        可上九天揽月，

        可下五洋捉鳖，

        谈笑凯歌还。

        世上无难事，

        只要肯登攀。

    \end{Large}

\end{center}

\vspace{8pt}
\clearpage

\section{画山水歌}

\beisong{泛读}

\begin{center}
    {\zihao{-0}\bfseries 画山水歌}

    \vspace{18pt}

    \begin{large}\xuwen
        〔吴融〕
    \end{large}

\end{center}

    \vspace{18pt}  % 起首

\begin{center}
    \begin{Large}

        良工善得丹青理，辄向茅茨画山水。

        地角移来方寸间，天涯写在笔锋里。

        日不落兮月长生，云片片兮水冷冷。

        经年蝴蝶飞不去，累岁桃花结不成。

        一块石，数株松，远又淡，近又浓。

        不出门庭三五步，观尽江山千万重。

    \end{Large}

\end{center}

\vspace{8pt}
\clearpage

\section{登幽州台歌}

\beisong{背诵}

\begin{center}
    {\zihao{-0}\bfseries 登幽州台歌}

    \vspace{18pt}

    \begin{large}\xuwen
        〔陈子昂〕
    \end{large}

\end{center}

    \vspace{18pt}  % 起首

\begin{center}
    \begin{Large}

        前不见古人，后不见来者。

        念天地之悠悠，独怆然而涕下。

    \end{Large}

\end{center}

\vspace{8pt}
\clearpage

\section{望岳}

\beisong{背诵}

\begin{center}
    {\zihao{-0}\bfseries 望岳}

    \vspace{18pt}

    \begin{large}\xuwen
        〔杜甫〕
    \end{large}

\end{center}

    \vspace{18pt}  % 起首

\begin{center}
    \begin{Large}

        岱宗夫如何？齐鲁青未了。

        造化钟神秀，阴阳割昏晓。

        荡胸生曾云，决眦入归鸟。

        会当凌绝顶，一览众山小。

    \end{Large}

\end{center}

\vspace{8pt}
\clearpage

\section{登飞来峰}

\beisong{背诵}

\begin{center}
    {\zihao{-0}\bfseries 登飞来峰}

    \vspace{18pt}

    \begin{large}\xuwen
        〔王安石〕
    \end{large}

\end{center}

    \vspace{18pt}  % 起首

\begin{center}
    \begin{Large}

        飞来山上千寻塔，闻说鸡鸣见日升。

        不畏浮云遮望眼，自缘身在最高层。

    \end{Large}

\end{center}

\vspace{8pt}
\clearpage

\section{游山西村}

\beisong{背诵}

\begin{center}
    {\zihao{-0}\bfseries 游山西村}

    \vspace{18pt}

    \begin{large}\xuwen
        〔陆游〕
    \end{large}

\end{center}

    \vspace{18pt}  % 起首

\begin{center}
    \begin{Large}

        莫笑农家腊酒浑，丰年留客足鸡豚。

        山重水复疑无路，柳暗花明又一村。

        萧鼓追随春社近，衣冠简朴古风存。

        从今若许闲乘月，拄杖无时夜叩门。

    \end{Large}

\end{center}

\vspace{8pt}
\clearpage

\section{己亥杂诗}

\beisong{泛读}

\begin{center}
    {\zihao{-0}\bfseries 己亥杂诗}

    \vspace{18pt}

    \begin{large}\xuwen
        〔龚自珍〕
    \end{large}

\end{center}

    \vspace{18pt}  % 起首

\begin{center}
    \begin{Large}

        浩荡离愁白日斜，吟鞭东指即天涯。

        落红不是无情物，化作春泥更护花。

    \end{Large}

\end{center}

\vspace{8pt}
\clearpage

\section{嫦娥}

\beisong{背诵}

\begin{center}
    {\zihao{-0}\bfseries 嫦娥}

    \vspace{18pt}

    \begin{large}\xuwen
        〔李商隐〕
    \end{large}

\end{center}

    \vspace{18pt}  % 起首

\begin{center}
    \begin{Large}

        云母屏风烛影深，长河渐落晓星沉。

        嫦娥应悔偷灵药，碧海青天夜夜心。

    \end{Large}

\end{center}

\vspace{8pt}
\clearpage

\section{浣溪沙}

\beisong{背诵}

\begin{center}
    {\zihao{-0}\bfseries 浣溪沙}

    \vspace{18pt}

    \begin{large}\xuwen
        〔苏轼〕
    \end{large}

\end{center}

    \vspace{12pt}\mbox{}

    \begin{large}\xuwen
        游蕲水清泉寺，寺临兰溪，溪水西流。
    \end{large}


    \vspace{8pt}

\begin{center}
    \begin{Large}

        山下兰芽短浸溪，松间沙路净无泥，

        潇潇暮雨子规啼。

        \vspace{18pt}  % 分阙

        谁道人生无再少？门前流水尚能西！

        休将白发唱黄鸡。

    \end{Large}

\end{center}

\vspace{8pt}
\clearpage

\section{春怨}

\beisong{泛读}

\begin{center}
    {\zihao{-0}\bfseries 春怨}

    \vspace{18pt}

    \begin{large}\xuwen
        〔金昌绪〕
    \end{large}

\end{center}

    \vspace{18pt}

\begin{center}
    \begin{Large}

        打起黄莺儿，莫教枝上啼。

        啼时惊妾梦，不得到辽西。

    \end{Large}

\end{center}

\vspace{8pt}
\clearpage

\section{贾生}

\beisong{背诵}

\begin{center}
    {\zihao{-0}\bfseries 贾生}

    \vspace{18pt}

    \begin{large}\xuwen
        〔李商隐〕
    \end{large}

\end{center}

    \vspace{18pt}

\begin{center}
    \begin{Large}

        宣室求贤访逐臣，贾生才调更无伦。

        可怜夜半虚前席，不问苍生问鬼神！

    \end{Large}

\end{center}

\vspace{8pt}
\clearpage

\chapter{第三级}
\clearpage

\section{沁园春·雪}

\beisong{背诵}

\begin{center}
    {\zihao{-0}\bfseries 沁园春·雪}

    \vspace{18pt}

    \begin{large}\xuwen
        〔毛泽东〕
    \end{large}

\end{center}

    \vspace{18pt}  % 起首

\begin{center}
    \begin{Large}

        北国风光，千里冰封，万里雪飘。

        望长城内外，惟余莽莽；

        大河上下，顿失滔滔。

        山舞银蛇，原驰蜡象，

        欲与天公试比高。

        须晴日，看红装素裹，分外妖娆。

        \vspace{18pt}  % 分阙

        江山如此多娇，

        引无数英雄竞折腰。

        惜秦皇汉武，略输文采；

        唐宗宋祖，稍逊风骚。

        一代天骄，成吉思汗，

        只识弯弓射大雕。

        俱往矣，数风流人物，还看今朝。

    \end{Large}

\end{center}

\vspace{8pt}
\clearpage

\section{桃源行}

\beisong{泛读}

\begin{center}
    {\zihao{-0}\bfseries 桃源行}

    \vspace{18pt}

    \begin{large}\xuwen
        〔王安石〕
    \end{large}

\end{center}

    \vspace{18pt}  % 起首

\begin{center}
    \begin{Large}

        望夷宫中鹿为马，秦人半死长城下。

        避时不独商山翁，亦有桃源种桃者。

        此来种桃经几春，采花食实枝为薪。

        儿孙生长与世隔，只有父子无君臣。

        渔郎漾舟迷远近，花间相见因相问。

        世上那知古有秦，山中岂料今为晋。

        闻道长安吹战尘，春风回首一沾巾。

        重华一去宁复得，天下纷纷经几秦。

    \end{Large}

\end{center}

\vspace{8pt}
\clearpage

\section{去矣行}

\beisong{泛读}

\begin{center}
    {\zihao{-0}\bfseries 去矣行}

    \vspace{18pt}

    \begin{large}\xuwen
        〔杜甫〕
    \end{large}

\end{center}

    \vspace{18pt}  % 起首

\begin{center}
    \begin{Large}

        君不见鞲上鹰，一饱则飞掣。

        焉能作堂上燕，衔泥附炎热。

        野人旷荡无腆颜，岂可久在王侯间。

        未试囊中餐玉法，明朝且入蓝田山。

    \end{Large}

\end{center}

\vspace{8pt}
\clearpage

\section{金陵酒肆留别}

\beisong{泛读}

\begin{center}
    {\zihao{-0}\bfseries 金陵酒肆留别}

    \vspace{18pt}

    \begin{large}\xuwen
        〔李白〕
    \end{large}

\end{center}

    \vspace{18pt}

\begin{center}
    \begin{Large}

        风吹柳花满店香，吴姬压酒唤客尝。

        金陵子弟来相送，欲行不行各尽觞。

        请君试问东流水，别意与之谁短长？

    \end{Large}

\end{center}

\vspace{8pt}
\clearpage

\section{子夜四时歌之秋歌}

\beisong{背诵}

\begin{center}
    {\zihao{-0}\bfseries 子夜四时歌之秋歌}

    \vspace{18pt}

    \begin{large}\xuwen
        〔李白〕
    \end{large}

\end{center}

    \vspace{18pt}

\begin{center}
    \begin{Large}

        长安一片月，万户捣衣声。

        秋风吹不尽，总是玉关情。

        何日平胡虏，良人罢远征？

    \end{Large}

\end{center}

\vspace{8pt}
\clearpage

\section{使至塞上}

\beisong{背诵}

\begin{center}
    {\zihao{-0}\bfseries 使至塞上}

    \vspace{18pt}

    \begin{large}\xuwen
        〔王维〕
    \end{large}

\end{center}

    \vspace{18pt}  % 起首

\begin{center}
    \begin{Large}

        单车欲问边，属国过居延，

        征蓬出汉塞，归雁入胡天，

        大漠孤烟直，长河落日圆。

        萧关逢侯骑，都护在燕然。

    \end{Large}

\end{center}

\vspace{8pt}
\clearpage

\section{渡荆门送别}

\beisong{背诵}

\begin{center}
    {\zihao{-0}\bfseries 渡荆门送别}

    \vspace{18pt}

    \begin{large}\xuwen
        〔李白〕
    \end{large}

\end{center}

    \vspace{18pt}  % 起首

\begin{center}
    \begin{Large}

        渡远荆门外，来从楚国游。

        山随平野尽，江入大荒流。

        月下飞天镜，云生结海楼。

        仍怜故乡水，万里送行舟。

    \end{Large}

\end{center}

\vspace{8pt}
\clearpage

\section{夜雨寄北}

\beisong{背诵}

\begin{center}
    {\zihao{-0}\bfseries 夜雨寄北}

    \vspace{18pt}

    \begin{large}\xuwen
        〔李商隐〕
    \end{large}

\end{center}

    \vspace{18pt}

\begin{center}
    \begin{Large}

        君问归期未有期，巴山夜雨涨秋池。

        何当共剪西窗烛，却话巴山夜雨时？

    \end{Large}

\end{center}

\vspace{8pt}
\clearpage

\section{钱塘湖春行}

\beisong{背诵}

\begin{center}
    {\zihao{-0}\bfseries 钱塘湖春行}

    \vspace{18pt}

    \begin{large}\xuwen
        〔白居易〕
    \end{large}

\end{center}

    \vspace{18pt}  % 起首

\begin{center}
    \begin{Large}

        孤山寺北贾亭西，水面初平云脚低。

        几处早莺争暖树，谁家新燕啄春泥。

        乱花渐欲迷人眼，浅草才能没马蹄。

        最爱湖东行不足，绿杨阴里白沙堤。

    \end{Large}

\end{center}

\vspace{8pt}
\clearpage

\section{饮酒}

\beisong{背诵}

\begin{center}
    {\zihao{-0}\bfseries 饮酒}

    \vspace{18pt}

    \begin{large}\xuwen
        〔陶渊明〕
    \end{large}

\end{center}

    \vspace{18pt}  % 起首

\begin{center}
    \begin{Large}

        结庐在人境，而无车马喧。

        问君何能尔，心远地自偏。

        采菊东篱下，悠然见南山。

        山气日夕佳，飞鸟相与还。

        此中有真意，欲辨已忘言。

    \end{Large}

\end{center}

\vspace{8pt}
\clearpage

\section{渡汉江}

\beisong{泛读}

\begin{center}
    {\zihao{-0}\bfseries 渡汉江}

    \vspace{18pt}

    \begin{large}\xuwen
        〔李频〕
    \end{large}

\end{center}

    \vspace{18pt}

\begin{center}
    \begin{Large}

        岭外音书绝，经冬复立春。

        近乡情更怯，不敢问来人。

    \end{Large}

\end{center}

\vspace{8pt}
\clearpage

\section{登金陵凤凰台}

\beisong{泛读}

\begin{center}
    {\zihao{-0}\bfseries 登金陵凤凰台}

    \vspace{18pt}

    \begin{large}\xuwen
        〔李白〕
    \end{large}

\end{center}

    \vspace{18pt}

\begin{center}
    \begin{Large}

        凤凰台上凤凰游，凤去台空江自流。

        吴宫花草埋幽径，晋代衣冠成古丘。

        三山半落青天外，二水中分白鹭洲。

        总为浮云能蔽日，长安不见使人愁。

    \end{Large}

\end{center}

\vspace{8pt}
\clearpage

\section{渔家傲}

\beisong{背诵}

\begin{center}
    {\zihao{-0}\bfseries 渔家傲}

    \vspace{18pt}

    \begin{large}\xuwen
        〔范仲淹〕
    \end{large}

\end{center}

    \vspace{18pt}  % 起首

\begin{center}
    \begin{Large}

        塞下秋来风景异，

        衡阳雁去无留意。

        四面边声连角起，

        千嶂里，长烟落日孤城闭。

        \vspace{18pt}  % 分阙

        浊酒一杯家万里，

        燕然未勒归无计。

        羌管悠悠霜满地，

        人不寐，将军白发征夫泪。

    \end{Large}

\end{center}

\vspace{8pt}
\clearpage

\section{旅夜书怀}

\beisong{背诵}

\begin{center}
    {\zihao{-0}\bfseries 旅夜书怀}

    \vspace{18pt}

    \begin{large}\xuwen
        〔杜甫〕
    \end{large}

\end{center}

    \vspace{18pt}  % 起首

\begin{center}
    \begin{Large}

        细草微风岸，危樯独夜舟。

        星垂平野阔，月涌大江流。

        名岂文章著，官应老病休。

        飘飘何所似，天地一沙鸥。

    \end{Large}

\end{center}

\vspace{8pt}
\clearpage

\chapter{第四级}
\clearpage

\section{忆秦娥·娄山关}

\beisong{背诵}

\begin{center}
    {\zihao{-0}\bfseries 忆秦娥·娄山关}

    \vspace{18pt}

    \begin{large}\xuwen
        〔毛泽东〕
    \end{large}

\end{center}

    \vspace{18pt}  % 起首

\begin{center}
    \begin{Large}

        西风烈，

        长空雁叫霜晨月。

        霜晨月，

        马蹄声碎，

        喇叭声咽。

        \vspace{18pt}  % 分阙

        雄关漫道真如铁，

        而今迈步从头越。

        从头越，

        苍山如海，

        残阳如血。

    \end{Large}

\end{center}

\vspace{8pt}
\clearpage

\section{观祈雨}

\beisong{泛读}

\begin{center}
    {\zihao{-0}\bfseries 观祈雨}

    \vspace{18pt}

    \begin{large}\xuwen
        〔李约〕
    \end{large}

\end{center}

    \vspace{18pt}  % 起首

\begin{center}
    \begin{Large}

        桑条无叶土生烟，箫管迎龙水庙前。

        朱门几处看歌舞，犹恐春阴咽管弦。

    \end{Large}

\end{center}

\vspace{8pt}
\clearpage

\section{莫相疑行}

\beisong{泛读}

\begin{center}
    {\zihao{-0}\bfseries 莫相疑行}

    \vspace{18pt}

    \begin{large}\xuwen
        〔杜甫〕
    \end{large}

\end{center}

    \vspace{18pt}  % 起首

\begin{center}
    \begin{Large}

        男儿生无所成头皓白，牙齿欲落真可惜。\phantom{男..}

        忆献三赋蓬莱宫，自怪一日声烜赫。

        集贤学士如堵墙，观我落笔中书堂。

        往时文采动人主，此日饥寒趋路旁。

        晚将末契托年少，当面输心背面笑。

        寄谢悠悠世上儿，不争好恶莫相疑。

    \end{Large}

\end{center}

\vspace{8pt}
\clearpage

\section{十五夜望月}

\beisong{泛读}

\begin{center}
    {\zihao{-0}\bfseries 十五夜望月}

    \vspace{18pt}

    \begin{large}\xuwen
        〔王建〕
    \end{large}

\end{center}

    \vspace{18pt}  % 起首

\begin{center}
    \begin{Large}

        中庭地白树栖鸦，冷露无声湿桂花。

        今夜月明人尽望，不知秋思落谁家？

    \end{Large}

\end{center}

\vspace{8pt}
\clearpage

\section{春望}

\beisong{背诵}

\begin{center}
    {\zihao{-0}\bfseries 春望}

    \vspace{18pt}

    \begin{large}\xuwen
        〔杜甫〕
    \end{large}

\end{center}

    \vspace{18pt}  % 起首

\begin{center}
    \begin{Large}

        国破山河在， 城春草木深。

        感时花溅泪， 恨别鸟惊心。

        烽火连三月， 家书抵万金。

        白头搔更短， 浑欲不胜簪。

    \end{Large}

\end{center}

\vspace{8pt}
\clearpage

\section{雁门太守行}

\beisong{背诵}

\begin{center}
    {\zihao{-0}\bfseries 雁门太守行}

    \vspace{18pt}

    \begin{large}\xuwen
        〔李贺〕
    \end{large}

\end{center}

    \vspace{18pt}  % 起首

\begin{center}
    \begin{Large}

        黑云压城城欲摧， 甲光向日金鳞开。

        角声满天秋色里， 塞上燕脂凝夜紫。

        半卷红旗临易水， 霜重鼓寒声不起。

        报君黄金台上意， 提携玉龙为君死。

    \end{Large}

\end{center}

\vspace{8pt}
\clearpage

\section{泊秦淮}

\beisong{背诵}

\begin{center}
    {\zihao{-0}\bfseries 泊秦淮}

    \vspace{18pt}

    \begin{large}\xuwen
        〔杜牧〕
    \end{large}

\end{center}

    \vspace{18pt}

\begin{center}
    \begin{Large}

        烟笼寒水月笼沙，夜泊秦淮近酒家。

        商女不知亡国恨，隔江犹唱后庭花。

    \end{Large}

\end{center}

\vspace{8pt}
\clearpage

\section{新嫁娘}

\beisong{泛读}

\begin{center}
    {\zihao{-0}\bfseries 新嫁娘}

    \vspace{18pt}

    \begin{large}\xuwen
        〔王建〕
    \end{large}

\end{center}

    \vspace{18pt}

\begin{center}
    \begin{Large}

        三日入厨下，洗手作羹汤。

        未谙姑食性，先遣小姑尝。

    \end{Large}

\end{center}

\vspace{8pt}
\clearpage

\section{汉江临眺}

\beisong{背诵}

\begin{center}
    {\zihao{-0}\bfseries 汉江临眺}

    \vspace{18pt}

    \begin{large}\xuwen
        〔王维〕
    \end{large}

\end{center}

    \vspace{18pt}

\begin{center}
    \begin{Large}

        楚塞三湘接，荆门九派通。

        江流天地外，山色有无中。

        郡邑浮前浦，波澜动远空。

        襄阳好风日，留醉与山翁。

    \end{Large}

\end{center}

\vspace{8pt}
\clearpage

\section{茅屋为秋风所破歌}

\beisong{背诵}

\begin{center}
    {\zihao{-0}\bfseries 茅屋为秋风所破歌}

    \vspace{18pt}

    \begin{large}\xuwen
        〔杜甫〕
    \end{large}

\end{center}

    \vspace{18pt}  % 起首

\begin{center}
    \begin{Large}

        八月秋高风怒号，卷我屋上三重茅。

        茅飞渡江洒江郊。

        高者挂罥长林梢，下者飘转沉塘坳。

        南村群童欺我老无力，忍能对面为盗贼。\hspace{24pt}

        公然抱茅入竹去。

        唇焦口燥呼不得，归来倚杖自叹息。

        俄顷风定云墨色，秋天漠漠向昏黑。

        布衾多年冷似铁，娇儿恶卧踏里裂。

        床头屋漏无干处，雨脚如麻未断绝。

        自经丧乱少睡眠，长夜沾湿何由彻！

        安得广厦千万间？
        
        大庇天下寒士俱欢颜，风雨不动安如山。\hspace{24pt}

        呜呼！
        
        何时眼前突兀见此屋，吾庐独破受冻死亦足！

    \end{Large}

\end{center}

\vspace{8pt}
\clearpage

\section{破阵子}

\beisong{背诵}

\begin{center}
    {\zihao{-0}\bfseries 破阵子}

    \vspace{18pt}

    \begin{large}\xuwen
        〔辛弃疾〕
    \end{large}

\end{center}

    \vspace{18pt}  % 起首

\begin{center}
    \begin{Large}

        醉里挑灯看剑，梦回吹角连营。

        八百里分麾下炙，五十弦翻塞外声。

        沙场秋点兵。

        \vspace{18pt}  % 分阙

        马作的卢飞快，弓如霹雳弦惊。

        了却君王天下事，赢得生前身后名。

        可怜白发生！

    \end{Large}

\end{center}

\vspace{8pt}
\clearpage

\section{过零丁洋}

\beisong{背诵}

\begin{center}
    {\zihao{-0}\bfseries 过零丁洋}

    \vspace{18pt}

    \begin{large}\xuwen
        〔文天祥〕
    \end{large}

\end{center}

    \vspace{18pt}  % 起首

\begin{center}
    \begin{Large}

        辛苦遭逢起一经，干戈寥落四周星。

        山河破碎风飘絮，身世浮沉雨打萍。

        惶恐滩头说惶恐，零丁洋里叹零丁。

        人生自古谁无死，留取丹心照汗青。

    \end{Large}

\end{center}

\vspace{8pt}
\clearpage

\section{越中览古}

\beisong{泛读}

\begin{center}
    {\zihao{-0}\bfseries 越中览古}

    \vspace{18pt}

    \begin{large}\xuwen
        〔李白〕
    \end{large}

\end{center}

    \vspace{18pt}  % 起首

\begin{center}
    \begin{Large}

        越王勾践破吴归，战士还家尽锦衣。

        宫女如花满春殿，只今惟有鹧鸪飞。

    \end{Large}

\end{center}

\vspace{8pt}
\clearpage

\chapter{第五级}
\clearpage

\section{水调歌头·游泳}

\beisong{精读}

\begin{center}
    {\zihao{-0}\bfseries 水调歌头·游泳}

    \vspace{18pt}

    \begin{large}\xuwen
        〔毛泽东〕
    \end{large}

\end{center}

    \vspace{18pt}  % 起首

\begin{center}
    \begin{Large}

        才饮长沙水，

        又食武昌鱼。

        万里长江横渡，

        极目楚天舒。

        不管风吹浪打，

        胜似闲庭信步，

        今日得宽馀。

        子在川上曰：

        逝者如斯夫！

        \vspace{18pt}  % 分阙

        风樯动，

        龟蛇静，

        起宏图。

        一桥飞架南北，

        天堑变通途。

        更立西江石壁，

        截断巫山云雨，

        高峡出平湖。

        神女应无恙，

        当惊世界殊。

    \end{Large}

\end{center}

\vspace{8pt}
\clearpage

\section{古长城吟}

\beisong{泛读}

\begin{center}
    {\zihao{-0}\bfseries 古长城吟}

    \vspace{18pt}

    \begin{large}\xuwen
        〔王翰〕
    \end{large}

\end{center}

    \vspace{18pt}  % 起首

\begin{center}
    \begin{Large}

        长安少年无远图，一生惟羡执金吾。

        麒麟前殿拜天子，走马西击长城胡。

        胡沙猎猎吹人面，汉虏相逢不相见。

        遥闻鼙鼓动地来，传道单于夜犹战。

        此时顾恩宁顾身，为君一行摧万人。

        壮士挥戈回白日，单于溅血染朱轮。

        归来饮马长城窟，长城道傍多白骨。

        问之耆老何代人，云是秦王筑城卒。

        黄昏塞北无人烟，鬼哭啾啾声沸天。

        无罪见诛功不赏，孤魂流落此城边。

        当昔秦王按剑起，诸侯膝行不敢视。

        富国强兵二十年，筑怨兴徭九千里。

        秦王筑城何太愚，天实亡秦非北胡。

        一朝祸起萧墙内，渭水咸阳不复都。

    \end{Large}

\end{center}

\vspace{8pt}
\clearpage

\section{木兰辞}

\beisong{背诵}

\begin{center}
    {\zihao{-0}\bfseries 木兰辞}

    \vspace{18pt}

    \begin{large}\xuwen
        〔民歌〕
    \end{large}

\end{center}

    \vspace{18pt}  % 起首

\begin{center}
    \begin{Large}

        唧唧复唧唧，木兰当户织。

        不闻机杼声，唯闻女叹息。

        问女何所思，问女何所忆。

        女亦无所思，女亦无所忆。

        昨夜见军帖，可汗大点兵，

        军书十二卷，卷卷有爷名。

        阿爷无大儿，木兰无长兄，

        愿为市鞍马，从此替爷征。

        东市买骏马，西市买鞍鞯，

        南市买辔头，北市买长鞭。

        旦辞爷娘去，暮宿黄河边，

        \hspace{36pt}不闻爷娘唤女声，但闻黄河流水鸣溅溅。

        旦辞黄河去，暮至黑山头，

        \hspace{36pt}不闻爷娘唤女声，但闻燕山胡骑鸣啾啾。

        万里赴戎机，关山度若飞。

        朔气传金柝，寒光照铁衣。

        将军百战死，壮士十年归。

        归来见天子，天子坐明堂。

        策勋十二转，赏赐百千强。

        \hspace{34pt}可汗问所欲，木兰不用尚书郎，

        愿驰千里足，送儿还故乡。

        爷娘闻女来，出郭相扶将；

        阿姊闻妹来，当户理红妆；

        \hspace{34pt}小弟闻姊来，磨刀霍霍向猪羊。

        开我东阁门，坐我西阁床。

        脱我战时袍，著我旧时裳。

        当窗理云鬓，对镜帖花黄。

        出门看火伴，火伴皆惊忙：

        \hspace{34pt}同行十二年，不知木兰是女郎。

        雄兔脚扑朔，雌兔眼迷离；

        \hspace{36pt}双兔傍地走，安能辨我是雄雌？

    \end{Large}

\end{center}

\vspace{8pt}
\clearpage

\section{渔歌子}

\beisong{泛读}

\begin{center}
    {\zihao{-0}\bfseries 渔歌子}

    \vspace{18pt}

    \begin{large}\xuwen
        〔张志和〕
    \end{large}

\end{center}

    \vspace{18pt}  % 起首

\begin{center}
    \begin{Large}

        西塞山前白鹭飞，桃花流水鳜鱼肥。

        青箬笠，绿蓑衣，斜风细雨不须归。

    \end{Large}

\end{center}

\vspace{8pt}
\clearpage


\section{问刘十九}

\beisong{背诵}

\begin{center}
    {\zihao{-0}\bfseries 问刘十九}

    \vspace{18pt}

    \begin{large}\xuwen
        〔白居易〕
    \end{large}

\end{center}

    \vspace{18pt}

\begin{center}
    \begin{Large}

        绿蚁新醅酒，红泥小火炉。

        晚来天欲雪，能饮一杯无？

    \end{Large}

\end{center}

\vspace{8pt}
\clearpage

\section{咏怀古迹}

\beisong{泛读}

\begin{center}
    {\zihao{-0}\bfseries 咏怀古迹}

    \vspace{18pt}

    \begin{large}\xuwen
        〔杜甫〕
    \end{large}

\end{center}

    \vspace{18pt}

\begin{center}
    \begin{Large}

        群山万壑赴荆门，生长明妃尚有村。

        一去紫台连朔漠，独留青冢向黄昏。

        画图省识春风面，环佩空归月下魂。

        千载琵琶作胡语，分明怨恨曲中论。

    \end{Large}

\end{center}

\vspace{8pt}
\clearpage

\section{诗经·蒹葭}

\beisong{背诵}

\begin{center}
    {\zihao{-0}\bfseries 诗经·蒹葭}

    \vspace{18pt}

    \begin{large}\xuwen
        〔诗经〕
    \end{large}

\end{center}

    \vspace{18pt}  % 起首

\begin{center}
    \begin{Large}

        蒹葭苍苍，白露为霜。

        所谓伊人，在水一方。

        溯洄从之，道阻且长。

        溯游从之，宛在水中央。

        \vspace{18pt}  % 分阙

        蒹葭萋萋，白露未晞。

        所谓伊人，在水之湄。

        溯洄从之，道阻且跻。

        溯游从之，宛在水中坻。

        \vspace{18pt}  % 分阙

        蒹葭采采，白露未已。

        所谓伊人，在水之涘。

        溯洄从之，道阻且右。

        溯游从之，宛在水中沚。

    \end{Large}

\end{center}

\vspace{8pt}
\clearpage

\section{江城子·密州出猎}

\beisong{背诵}

\begin{center}
    {\zihao{-0}\bfseries 江城子·密州出猎}

    \vspace{18pt}

    \begin{large}\xuwen
        〔苏轼〕
    \end{large}

\end{center}

    \vspace{18pt}  % 起首

\begin{center}
    \begin{Large}

        老夫聊发少年狂，

        左牵黄，右擎苍，

        锦帽貂裘，千骑卷平冈。

        为报倾城随太守，

        亲射虎，看孙郎。

        \vspace{18pt}  % 分阙

        酒酣胸胆尚开张，

        鬓微霜，又何妨？

        持节云中，何日遣冯唐。

        会挽雕弓如满月，

        西北望，射天狼。

    \end{Large}

\end{center}

\vspace{8pt}
\clearpage

\section{杂诗}

\beisong{泛读}

\begin{center}
    {\zihao{-0}\bfseries 杂诗}

    \vspace{18pt}

    \begin{large}\xuwen
        〔无名氏〕
    \end{large}

\end{center}

    \vspace{18pt}

\begin{center}
    \begin{Large}

        尽寒食雨草萋萋，著麦苗风柳映堤。

        等是有家归未得，杜鹃休向耳边啼。

    \end{Large}

\end{center}

\vspace{8pt}
\clearpage

\section{十五从军征}

\beisong{泛读}

\begin{center}
    {\zihao{-0}\bfseries 十五从军征}

    \vspace{18pt}

    \begin{large}\xuwen
        〔乐府〕
    \end{large}

\end{center}

    \vspace{18pt}  % 起首

\begin{center}
    \begin{Large}

        十五从军征，八十始得归。

        道逢乡里人：家中有阿谁？

        遥看是君家，松柏冢累累。

        兔从狗窦入，雉从梁上飞。

        中庭生旅谷，井上生旅葵。

        舂谷持作饭，采葵持作羹。

        羹饭一时熟，不知贻阿谁！

        出门东向看，泪落沾我衣。

    \end{Large}

\end{center}

\vspace{8pt}
\clearpage

\section{白雪歌送武判官归京}

\beisong{精读}

\begin{center}
    {\zihao{-0}\bfseries 白雪歌送武判官归京}

    \vspace{18pt}

    \begin{large}\xuwen
        〔岑参〕
    \end{large}

\end{center}

    \vspace{18pt}  % 起首

\begin{center}
    \begin{Large}

        北风卷地白草折，胡天八月即飞雪。

        忽如一夜春风来，千树万树梨花开。

        散入珠帘湿罗幕，狐裘不暖锦衾薄。

        将军角弓不得控，都护铁衣冷难着。

        瀚海阑干百丈冰，愁云惨淡万里凝。

        中军置酒饮归客，胡琴琵琶与羌笛。

        纷纷暮雪下辕门，风掣红旗冻不翻。

        轮台东门送君去，去时雪满天山路。

        山回路转不见君，雪上空留马行处。

    \end{Large}

\end{center}

\vspace{8pt}
\clearpage

\section{南乡子}

\beisong{背诵}

\begin{center}
    {\zihao{-0}\bfseries 南乡子}

    \vspace{18pt}

    \begin{large}\xuwen
        〔辛弃疾〕
    \end{large}

\end{center}

    \vspace{18pt}  % 起首

\begin{center}
    \begin{Large}

        何处望神州？满眼风光北固楼。

        千古兴亡多少事？悠悠。

        不尽长江滚滚流。

        \vspace{18pt}  % 分阙

        年少万兜鍪，坐断东南战未休。

        天下英雄谁敌手？曹刘。

        生子当如孙仲谋。

    \end{Large}

\end{center}

\vspace{8pt}
\clearpage

\chapter{第六级}
\clearpage

\section{水调歌头}

\beisong{背诵}

\begin{center}
    {\zihao{-0}\bfseries 水调歌头}

    \vspace{18pt}

    \begin{large}\xuwen
        〔苏轼〕
    \end{large}

\end{center}

    \vspace{18pt}  % 起首

\begin{center}
    \begin{Large}

        明月几时有？把酒问青天。

        不知天上宫阙，今夕是何年。

        我欲乘风归去，又恐琼楼玉宇，高处不胜寒。

        起舞弄清影，何似在人间。

        \vspace{18pt}  % 分阙

        转朱阁，低绮户，照无眠。

        不应有恨，何事长向别时圆？

        人有悲欢离合，月有阴晴圆缺，此事古难全。

        但愿人长久，千里共婵娟。

    \end{Large}

\end{center}

\vspace{8pt}
\clearpage

\section{赤壁}

\beisong{背诵}

\begin{center}
    {\zihao{-0}\bfseries 赤壁}

    \vspace{18pt}

    \begin{large}\xuwen
        〔杜牧〕
    \end{large}

\end{center}

    \vspace{18pt}  % 起首

\begin{center}
    \begin{Large}

        折戟沉沙铁未销，自将磨洗认前朝。

        东风不与周郎便，铜雀春深锁二乔。

    \end{Large}

\end{center}

\vspace{8pt}
\clearpage

\section{诗经·关雎}

\beisong{背诵}

\begin{center}
    {\zihao{-0}\bfseries 诗经·关雎}

    \vspace{18pt}

    \begin{large}\xuwen
        〔诗经〕
    \end{large}

\end{center}

    \vspace{18pt}  % 起首

\begin{center}
    \begin{Large}

        关关雎鸠，在河之洲。

        窈窕淑女，君子好逑。

        参差荇菜，左右流之。

        窈窕淑女，寤寐求之。

        求之不得，寤寐思服。

        悠哉悠哉，辗转反侧。

        参差荇菜，左右采之。

        窈窕淑女，琴瑟友之。

        参差荇菜，左右芼之。

        窈窕淑女，钟鼓乐之。

    \end{Large}

\end{center}

\vspace{8pt}
\clearpage

\section{酬乐天扬州初逢席上见赠}

\beisong{背诵}

\begin{center}
    {\zihao{-0}\bfseries 酬乐天扬州初逢席上见赠}

    \vspace{18pt}

    \begin{large}\xuwen
        〔刘禹锡〕
    \end{large}

\end{center}

    \vspace{18pt}  % 起首

\begin{center}
    \begin{Large}

        巴山楚水凄凉地，二十三年弃置身。

        怀旧空吟闻笛赋，到乡翻似烂柯人。

        沉舟侧畔千帆过，病树前头万木春。

        今日听君歌一曲，暂凭杯酒长精神。

    \end{Large}

\end{center}

\vspace{8pt}
\clearpage

\section{哥舒歌}

\beisong{泛读}

\begin{center}
    {\zihao{-0}\bfseries 哥舒歌}

    \vspace{18pt}

    \begin{large}\xuwen
        〔西鄙人〕
    \end{large}

\end{center}

    \vspace{18pt}

\begin{center}
    \begin{Large}

        北斗七星高，哥舒夜带刀。

        至今窥牧马，不敢过临洮。

    \end{Large}

\end{center}

\vspace{8pt}
\clearpage

\section{金缕衣}

\beisong{泛读}

\begin{center}
    {\zihao{-0}\bfseries 金缕衣}

    \vspace{18pt}

    \begin{large}\xuwen
        〔杜秋娘〕
    \end{large}

\end{center}

    \vspace{18pt}

\begin{center}
    \begin{Large}

        劝君莫惜金缕衣，劝君惜取少年时。

        花开堪折直须折，莫待无花空折枝！

    \end{Large}

\end{center}

\vspace{8pt}
\clearpage

\section{江南逢李龟年}

\beisong{泛读}

\begin{center}
    {\zihao{-0}\bfseries 江南逢李龟年}

    \vspace{18pt}

    \begin{large}\xuwen
        〔杜甫〕
    \end{large}

\end{center}

    \vspace{18pt}

\begin{center}
    \begin{Large}

        岐王宅里寻常见，崔九堂前几度闻。

        正是江南好风景，落花时节又逢君。

    \end{Large}

\end{center}

\vspace{8pt}
\clearpage

\section{满江红}

\beisong{背诵}

\begin{center}
    {\zihao{-0}\bfseries 满江红}

    \vspace{18pt}

    \begin{large}\xuwen
        〔岳飞〕
    \end{large}

\end{center}

    \vspace{18pt}  % 起首

\begin{center}
    \begin{Large}

        怒发冲冠，凭阑处、潇潇雨歇。

        抬望眼，仰天长啸，壮怀激烈。
        
        三十功名尘与土，八千里路云和月。
        
        莫等闲，白了少年头，空悲切。

        \vspace{18pt}  % 分阙

        靖康耻，犹未雪；

        臣子恨，何时灭？

        驾长车，踏破贺兰山缺。

        壮志饥餐胡虏肉，笑谈渴饮匈奴血。

        待从头，收拾旧山河，朝天阙。

    \end{Large}

\end{center}

\vspace{8pt}
\clearpage

\section{山坡羊·潼关怀古}

\beisong{背诵}

\begin{center}
    {\zihao{-0}\bfseries 山坡羊·潼关怀古}

    \vspace{18pt}

    \begin{large}\xuwen
        〔张养浩〕
    \end{large}

\end{center}

    \vspace{18pt}  % 起首

\begin{center}
    \begin{Large}

        峰峦如聚，波涛如怒，

        山河表里潼关路。

        望西都，意踌躇。

        \vspace{18pt}  % 分阙

        伤心秦汉经行处，

        宫阙万间都做了土。

        兴，百姓苦；亡，百姓苦。

    \end{Large}

\end{center}

\vspace{8pt}
\clearpage

\section{诗经·桃夭}

\beisong{泛读}

\begin{center}
    {\zihao{-0}\bfseries 诗经·桃夭}

    \vspace{18pt}

    \begin{large}\xuwen
        〔诗经〕
    \end{large}

\end{center}

    \vspace{18pt}  % 起首

\begin{center}
    \begin{Large}

        桃之夭夭，灼灼其华。

        之子于归，宜其室家。

        \vspace{18pt}  % 分阙

        桃之夭夭，有蕡其实。

        之子于归，宜其家室。

        \vspace{18pt}  % 分阙

        桃之夭夭，其叶蓁蓁。

        之子于归，宜其家人。

    \end{Large}

\end{center}

\vspace{8pt}
\clearpage

\section{汾阴行}

\beisong{泛读}

\begin{center}
    {\zihao{-0}\bfseries 汾阴行}

    \vspace{18pt}

    \begin{large}\xuwen
        〔李峤〕
    \end{large}

\end{center}

    \vspace{18pt}  % 起首

\begin{center}
    \begin{Large}

        君不见昔日西京全盛时，汾阴后土亲祭祀。\hspace{40pt}

        斋宫宿寝设储供，撞钟鸣鼓树羽旂。

        汉家五叶才且雄，宾延万灵朝九戎。

        柏梁赋诗高宴罢，诏书法驾幸河东。

        河东太守亲扫除，奉迎至尊导鸾舆。

        五营夹道列容卫，三河纵观空里闾。

        回旌驻跸降灵场，焚香奠醑邀百祥。

        金鼎发色正焜煌，灵祇炜烨摅景光。

        埋玉陈牲礼神毕，举麾上马乘舆出。

        彼汾之曲嘉可游，木兰为楫桂为舟。

        櫂歌微吟彩鹢浮，箫鼓哀鸣白云秋。

        欢娱宴洽赐群后，家家复除户牛酒。

        声明动天乐无有，千秋万岁南山寿。

        自从天子向秦关，玉辇金车不复还。

        珠帘羽扇长寂寞，鼎湖龙髯安可攀。

        千龄人事一朝空，四海为家此路穷。

        豪雄意气今何在，坛场宫馆尽蒿蓬。

        路逢故老长叹息，世事回环不可测。

        昔时青楼对歌舞，今日黄埃聚荆棘。

        山川满目泪沾衣，富贵荣华能几时？

        不见只今汾水上，唯有年年秋雁飞。

    \end{Large}

\end{center}

\vspace{8pt}
\clearpage

\section{梁甫吟}

\beisong{泛读}

\begin{center}
    {\zihao{-0}\bfseries 梁甫吟}

    \vspace{18pt}

    \begin{large}\xuwen
        〔佚名〕
    \end{large}

\end{center}

    \vspace{18pt}  % 起首

\begin{center}
    \begin{Large}

        步出齐城门，遥望荡阴里。

        里中有三坟，累累正相似。

        问是谁家墓，田疆古冶子。

        力能排南山，文能绝地理。

        一朝被谗言，二桃杀三士。

        谁能为此谋，相国齐晏子。

    \end{Large}

\end{center}

\vspace{8pt}
\clearpage

\section{七夕歌}

\beisong{泛读}

\begin{center}
    {\zihao{-0}\bfseries 七夕歌}

    \vspace{18pt}

    \begin{large}\xuwen
        〔张文潜〕
    \end{large}

\end{center}

    \vspace{18pt}  % 起首

\begin{center}
    \begin{Large}

        人间一叶梧桐飘，蓐收行秋回斗杓。

        神官召集役灵鹊，直渡天河横作桥。

        河东美人天帝子，机杼年年劳玉指。

        织成云雾紫绡衣，辛苦无欢容不理。

        帝怜独居无与娱，河西嫁与牵牛夫。

        自从嫁后废织纴，绿鬓云鬟朝暮梳。

        贪欢不归天帝怒，谪归却踏来时路。

        但令一岁一相逢，七月七夕河边渡。

        别多会少知奈何，却悔从前恩爱多。

        匆匆离恨说不尽，烛龙已驾随羲和。

        河边灵官晓催发，令严不管轻离别。

        空将泪作雨滂沱，泪痕有尽愁无歇。

        寄言织女君休叹，天地无穷会相见。

        犹胜姮娥不嫁人，夜夜孤眠广寒殿。

    \end{Large}

\end{center}

\vspace{8pt}
\clearpage


\chapter{第七级}
\clearpage

\section{沁园春·长沙}

\beisong{背诵}

\begin{center}
    {\zihao{-0}\bfseries 沁园春·长沙}

    \vspace{18pt}

    \begin{large}\xuwen
        〔毛泽东〕
    \end{large}

\end{center}

    \vspace{18pt}  % 起首

\begin{center}
    \begin{Large}

        独立寒秋，湘江北去，橘子洲头。

        看万山红遍，层林尽染；

        漫江碧透，百舸争流。

        鹰击长空，鱼翔浅底，

        万类霜天竞自由。

        怅寥廓，问苍茫大地：谁主沉浮？

        \vspace{18pt}  % 分阙

        携来百侣曾游。

        忆往昔峥嵘岁月稠。

        恰同学少年，风华正茂；

        书生意气，挥斥方遒。

        指点江山，激扬文字，

        粪土当年万户侯。

        曾记否，到中流击水，浪遏飞舟？

    \end{Large}

\end{center}

\vspace{8pt}
\clearpage

\section{草书歌行}

\beisong{泛读}

\begin{center}
    {\zihao{-0}\bfseries 草书歌行}

    \vspace{18pt}

    \begin{large}\xuwen
        〔李白〕
    \end{large}

\end{center}

    \vspace{18pt}  % 起首

\begin{center}
    \begin{Large}

        少年上人号怀素，草书天下称独步。

        墨池飞出北溟鱼，笔锋杀尽中山兔。

        八月九月天气凉，酒徒词客满高堂。

        笺麻素绢排数箱，宣州石砚墨色光。

        吾师醉后倚绳床，须臾扫尽数千张。

        飘风骤雨惊飒飒，落花飞雪何茫茫。

        起来向壁不停手，一行数字大如斗。

        怳怳如闻神鬼惊，时时只见龙蛇走。

        左盘右蹙如惊电，状同楚汉相攻战。

        湖南七郡凡几家，家家屏障书题遍。

        王逸少，张伯英，古来几许浪得名。

        张颠老死不足数，我师此义不师古。

        \hspace{52pt}古来万事贵天生，何必要公孙大娘浑脱舞。

    \end{Large}

\end{center}

\vspace{8pt}
\clearpage

\section{宣州谢朓楼饯别校书叔云}

\beisong{背诵}

\begin{center}
    {\zihao{-0}\bfseries 宣州谢朓楼饯别校书叔云}

    \vspace{18pt}

    \begin{large}\xuwen
        〔李白〕
    \end{large}

\end{center}

    \vspace{18pt}

\begin{center}
    \begin{Large}

        弃我去者，昨日之日不可留。

        乱我心者，今日之日多烦忧！

        长风万里送秋雁，对此可以酣高楼。

        蓬莱文章建安骨，中间小谢又清发。

        俱怀逸兴壮思飞，欲上青天览明月。

        抽刀断水水更流，举杯销愁愁更愁。

        人生在世不称意，明朝散发弄扁舟。

    \end{Large}

\end{center}

\vspace{8pt}
\clearpage

\section{卖炭翁}

\beisong{精读}

\begin{center}
    {\zihao{-0}\bfseries 卖炭翁}

    \vspace{18pt}

    \begin{large}\xuwen
        〔白居易〕
    \end{large}

\end{center}

    \vspace{18pt}  % 起首

\begin{center}
    \begin{Large}

        卖炭翁，卖炭翁，伐薪烧炭南山中。

        满面尘灰烟火色，两鬓苍苍十指黑。

        卖炭得钱何所营？身上衣裳口中食。

        可怜身上衣正单，心忧炭贱愿天寒。

        夜来城外一尺雪，晓驾炭车辗冰辙。

        牛困人饥日已高，市南门外泥中歇。

        翩翩两骑来是谁？黄衣使者白衫儿。

        手把文书口称敕，回车叱牛牵向北。

        一车炭，千余斤，宫使驱将惜不得。

        半匹红绡一丈绫，系向牛头充炭直。

    \end{Large}

\end{center}

\vspace{8pt}
\clearpage

\section{送上人}

\beisong{泛读}

\begin{center}
    {\zihao{-0}\bfseries 送上人}

    \vspace{18pt}

    \begin{large}\xuwen
        〔刘长卿〕
    \end{large}

\end{center}

    \vspace{18pt}

\begin{center}
    \begin{Large}

        孤云将野鹤，岂向人间住！

        莫买沃洲山，时人已知处。

    \end{Large}

\end{center}

\vspace{8pt}
\clearpage

\section{行路难}

\beisong{背诵}

\begin{center}
    {\zihao{-0}\bfseries 行路难}

    \vspace{18pt}

    \begin{large}\xuwen
        〔李白〕
    \end{large}

\end{center}

    \vspace{18pt}

\begin{center}
    \begin{Large}

        金樽清酒斗十千，玉盘珍羞值万钱。

        停杯投箸不能食，拔剑四顾心茫然。

        欲渡黄河冰塞川，将登太行雪满山。

        闲来垂钓碧溪上，忽复乘舟梦日边。

        行路难，行路难！多歧路，今安在？

        长风破浪会有时，直挂云帆济沧海。

    \end{Large}

\end{center}

\vspace{8pt}
\clearpage

\section{诗经·无衣}

\beisong{精读}

\begin{center}
    {\zihao{-0}\bfseries 诗经·无衣}

    \vspace{18pt}

    \begin{large}\xuwen
        〔诗经〕
    \end{large}

\end{center}

    \vspace{18pt}  % 起首

\begin{center}
    \begin{Large}

        岂曰无衣？与子同袍。

        王于兴师，修我戈矛。

        与子同仇！\hspace{78pt}

        \vspace{18pt}  % 分阙

        岂曰无衣？与子同泽。

        王于兴师，修我矛戟。

        与子偕作！\hspace{78pt}

        \vspace{18pt}  % 分阙

        岂曰无衣？与子同裳。

        王于兴师，修我甲兵。

        与子偕行！\hspace{78pt}

    \end{Large}

\end{center}

\vspace{8pt}
\clearpage

\section{涉江采芙蓉}

\beisong{泛读}

\begin{center}
    {\zihao{-0}\bfseries 涉江采芙蓉}

    \vspace{18pt}

    \begin{large}\xuwen
        〔无名氏〕
    \end{large}

\end{center}

    \vspace{12pt}\mbox{}

\begin{large}\xuwen
    取自古诗十九首。古诗十九首，作者不详，一般认为是东汉后期所作。最早由南朝梁萧统选入《昭明文选》后定名，习惯以首句为标题。
\end{large}

    \vspace{12pt}

\begin{center}
    \begin{Large}

        涉江采芙蓉， 兰泽多芳草。

        采之欲遗谁， 所思在远道。

        还顾望旧乡， 长路漫浩浩。

        同心而离居， 忧伤以终老。

    \end{Large}

\end{center}

\vspace{8pt}
\clearpage

\section{短歌行}

\beisong{精读}

\begin{center}
    {\zihao{-0}\bfseries 短歌行}

    \vspace{18pt}

    \begin{large}\xuwen
        〔曹操〕
    \end{large}

\end{center}

    \vspace{12pt}

\begin{center}
    \begin{Large}

        对酒当歌，人生几何！

        譬如朝露，去日苦多。

        慨当以慷，忧思难忘。

        何以解忧？唯有杜康。

        \vspace{18pt}  % 分阙

        青青子衿，悠悠我心。

        但为君故，沉吟至今。

        呦呦鹿鸣，食野之苹。

        我有嘉宾，鼓瑟吹笙。

        \vspace{18pt}  % 分阙

        明明如月，何时可掇？

        忧从中来，不可断绝。

        越陌度阡，枉用相存。

        契阔谈讌，心念旧恩。

        \vspace{18pt}  % 分阙

        月明星稀，乌鹊南飞，

        绕树三匝，何枝可依？

        山不厌高，海不厌深。

        周公吐哺，天下归心。

    \end{Large}

\end{center}

\vspace{8pt}
\clearpage

\section{陇西行}

\beisong{背诵}

\begin{center}
    {\zihao{-0}\bfseries 陇西行}

    \vspace{18pt}

    \begin{large}\xuwen
        〔陈陶〕
    \end{large}

\end{center}

    \vspace{18pt}

\begin{center}
    \begin{Large}

        誓扫匈奴不顾身，五千貂锦丧胡尘。

        可怜无定河边骨，犹是深闺梦里人！

    \end{Large}

\end{center}

\vspace{8pt}
\clearpage

\section{山居秋暝}

\beisong{背诵}

\begin{center}
    {\zihao{-0}\bfseries 山居秋暝}

    \vspace{18pt}

    \begin{large}\xuwen
        〔王维〕
    \end{large}

\end{center}

    \vspace{18pt}  % 起首

\begin{center}
    \begin{Large}

        空山新雨后，天气晚来秋。

        明月松间照，清泉石上流。

        竹喧归浣女，莲动下渔舟。

        随意春芳歇，王孙自可留。

    \end{Large}

\end{center}

\vspace{8pt}
\clearpage

\section{青玉案·元夕}

\beisong{背诵}

\begin{center}
    {\zihao{-0}\bfseries 青玉案·元夕}

    \vspace{18pt}

    \begin{large}\xuwen
        〔辛弃疾〕
    \end{large}

\end{center}

    \vspace{18pt}  % 起首

\begin{center}
    \begin{Large}

        东风夜放花千树。更吹落、星如雨。

        宝马雕车香满路。

        凤箫声动，玉壶光转，一夜鱼龙舞。

        \vspace{18pt}  % 分阙

        蛾儿雪柳黄金缕。笑语盈盈暗香去。

        众里寻他千百度。

        蓦然回首，那人却在，灯火阑珊处。

    \end{Large}

\end{center}

\vspace{8pt}
\clearpage

\section{石壕吏}

\beisong{背诵}

\begin{center}
    {\zihao{-0}\bfseries 石壕吏}

    \vspace{18pt}

    \begin{large}\xuwen
        〔杜甫〕
    \end{large}

\end{center}

    \vspace{18pt}  % 起首

\begin{center}
    \begin{Large}

        暮投石壕村，有吏夜捉人。

        老翁逾墙走，老妇出门看。

        吏呼一何怒，妇啼一何苦。

        听妇前致词：三男邺城戍，

        一男附书至，二男新战死。

        存者且偷生，死者长已矣。

        室中更无人，惟有乳下孙。

        有孙母未去，出入无完裙。

        老妪力虽衰，请从吏夜归。

        急应河阳役，犹得备晨炊。

        夜久语声绝，如闻泣幽咽。

        天明登前途，独与老翁别。

    \end{Large}

\end{center}

\vspace{8pt}
\clearpage

\chapter{第八级}
\clearpage

\section{琵琶行}

\beisong{背诵}

\begin{center}
    {\zihao{-0}\bfseries 琵琶行}

    \vspace{18pt}

    \begin{large}\xuwen
        〔白居易〕
    \end{large}
    
\end{center}

    \vspace{12pt}\mbox{}

    \begin{large}\xuwen
        元和十年，予左迁九江郡司马。明年秋，送客湓浦口，闻舟中夜弹琵琶者，听其音，铮铮然有京都声。问其人，本长安倡女，尝学琵琶于穆、曹二善才，年长色衰，委身为贾人妇。遂命酒，使快弹数曲。曲罢悯然，自叙少小时欢乐事，今漂沦憔悴，转徙于江湖间。予出官二年，恬然自安，感斯人言，是夕始觉有迁谪意。因为长句，歌以赠之，凡六百一十六言，命曰《琵琶行》。
    \end{large}


    \vspace{12pt}

\begin{center}
    \begin{Large}

        浔阳江头夜送客，枫叶荻花秋瑟瑟。

        主人下马客在船，举酒欲饮无管弦。

        醉不成欢惨将别，别时茫茫江浸月。

        忽闻水上琵琶声，主人忘归客不发。

        寻声暗问弹者谁，琵琶声停欲语迟。

        移船相近邀相见，添酒回灯重开宴。

        千呼万唤始出来，犹抱琵琶半遮面。

        转轴拨弦三两声，未成曲调先有情。

        弦弦掩抑声声思，似诉平生不得志。

        低眉信手续续弹，说尽心中无限事。

       \hspace{58pt} 轻拢慢捻抹复挑，初为《霓裳》后《六幺》。

        大弦嘈嘈如急雨，小弦切切如私语。

        嘈嘈切切错杂弹，大珠小珠落玉盘。

        间关莺语花底滑，幽咽泉流冰下难。

        冰泉冷涩弦凝绝，凝绝不通声暂歇。

        别有幽愁暗恨生，此时无声胜有声。

        银瓶乍破水浆迸，铁骑突出刀枪鸣。

        曲终收拨当心画，四弦一声如裂帛。

        东船西舫悄无言，唯见江心秋月白。

        沉吟放拨插弦中，整顿衣裳起敛容。

        自言本是京城女，家在虾蟆陵下住。

        十三学得琵琶成，名属教坊第一部。

        曲罢曾教善才服，妆成每被秋娘妒。

        五陵年少争缠头，一曲红绡不知数。

        钿头银篦击节碎，血色罗裙翻酒污。

        今年欢笑复明年，秋月春风等闲度。

        弟走从军阿姨死，暮去朝来颜色故。

        门前冷落鞍马稀，老大嫁作商人妇。

        商人重利轻别离，前月浮梁买茶去。

        去来江口守空船，绕船月明江水寒。

        夜深忽梦少年事，梦啼妆泪红阑干。

        我闻琵琶已叹息，又闻此语重唧唧。

        同是天涯沦落人，相逢何必曾相识！

        我从去年辞帝京，谪居卧病浔阳城。

        浔阳地僻无音乐，终岁不闻丝竹声。

        住近湓江地低湿，黄芦苦竹绕宅生。

        其间旦暮闻何物？杜鹃啼血猿哀鸣。

        春江花朝秋月夜，往往取酒还独倾。

        岂无山歌与村笛，呕哑嘲哳难为听。

        今夜闻君琵琶语，如听仙乐耳暂明。

        \hspace{26pt}莫辞更坐弹一曲，为君翻作《琵琶行》。

        感我此言良久立，却坐促弦弦转急。

        凄凄不似向前声，满座重闻皆掩泣。

        座中泣下谁最多？江州司马青衫湿。

    \end{Large}

\end{center}

\vspace{8pt}
\clearpage

\section{百舌吟}

\beisong{泛读}

\begin{center}
    {\zihao{-0}\bfseries 百舌吟}

    \vspace{18pt}

    \begin{large}\xuwen
        〔刘禹锡〕
    \end{large}

\end{center}

    \vspace{18pt}  % 起首

\begin{center}
    \begin{Large}

        晓星寥落春云低，初闻百舌间关啼。

        花树满空迷处所，摇动繁英坠红雨。

        笙簧百啭音韵多，黄鹂吞声燕无语。

        东方朝日迟迟升，迎风弄景如自矜。

        数声不尽又飞去，何许相逢绿杨路。

        绵蛮宛转似娱人，一心百舌何纷纷。

        酡颜侠少停歌听，坠珥妖姬和睡闻。

        可怜光景何时尽，谁能低回避鹰隼。

        廷尉张罗自不关，潘郎挟弹无情损。

        天生羽族尔何微，舌端万变乘春晖。

        南方朱鸟一朝见，索寞无言蒿下飞。

    \end{Large}

\end{center}

\vspace{8pt}
\clearpage

\section{观沧海}

\beisong{背诵}

\begin{center}
    {\zihao{-0}\bfseries 观沧海}

    \vspace{18pt}

    \begin{large}\xuwen
        〔曹操〕
    \end{large}

\end{center}

    \vspace{18pt}  % 起首

\begin{center}
    \begin{Large}

        东临碣石，以观沧海。

        水何澹澹，山岛竦峙。

        树木丛生，百草丰茂。

        秋风萧瑟，洪波涌起。

        日月之行，若出其中；

        星汉灿烂，若出其里。

        幸甚至哉，歌以咏志。

    \end{Large}

\end{center}

\vspace{8pt}
\clearpage

\section{离骚}

\beisong{精读}

\begin{center}
    {\zihao{-0}\bfseries 离骚}

    \vspace{18pt}

    \begin{large}\xuwen
        〔屈原〕
    \end{large}

\end{center}

    \vspace{18pt}  % 起首

\begin{center}
    \begin{Large}

        帝高阳之苗裔兮， 朕皇考曰伯庸。

        摄提贞于孟陬兮， 惟庚寅吾以降。

        皇览揆余初度兮， 肇锡余以嘉名。

        名余曰正则兮， 字余曰灵均。

        纷吾既有此内美兮， 又重之以修能。\hspace{5pt}

        扈江离与辟芷兮， 纫秋兰以为佩。

        \hspace{18pt}汨余若将不及兮， 恐年岁之不吾与。

        朝搴阰之木兰兮， 夕揽洲之宿莽。

        日月忽其不淹兮， 春与秋其代序。

        惟草木之零落兮， 恐美人之迟暮。

        不抚壮而弃秽兮， 何不改乎此度？

        乘骐骥以驰骋兮， 来吾导夫先路！

    \end{Large}

\end{center}

\vspace{8pt}
\clearpage

\section{无题}

\beisong{背诵}

\begin{center}
    {\zihao{-0}\bfseries 无题}

    \vspace{18pt}

    \begin{large}\xuwen
        〔李商隐〕
    \end{large}

\end{center}

    \vspace{18pt}

\begin{center}
    \begin{Large}

        相见时难别亦难，东风无力百花残。

        春蚕到死丝方尽，蜡炬成灰泪始干。

        晓镜但愁云鬓改，夜吟应觉月光寒。

        蓬莱此去无多路，青鸟殷勤为探看。

    \end{Large}

\end{center}

\vspace{8pt}
\clearpage

\section{无题}

\beisong{泛读}

\begin{center}
    {\zihao{-0}\bfseries 无题}

    \vspace{18pt}

    \begin{large}\xuwen
        〔李商隐〕
    \end{large}

\end{center}

    \vspace{18pt}

\begin{center}
    \begin{Large}

        来是空言去绝踪，月斜楼上五更钟。

        梦为远别啼难唤，书被催成墨未浓。

        蜡照半笼金翡翠，麝熏微度绣芙蓉。

        刘郎已恨蓬山远，更隔蓬山一万重。

    \end{Large}

\end{center}

\vspace{8pt}
\clearpage

\section{无题}

\beisong{泛读}

\begin{center}
    {\zihao{-0}\bfseries 无题}

    \vspace{18pt}

    \begin{large}\xuwen
        〔李商隐〕
    \end{large}

\end{center}

    \vspace{18pt}

\begin{center}
    \begin{Large}

        昨夜星辰昨夜风，画楼西畔桂堂东。

        身无彩凤双飞翼，心有灵犀一点通。

        隔座送钩春酒暖，分曹射覆蜡灯红。

        嗟余听鼓应官去，走马兰台类转蓬。

    \end{Large}

\end{center}

\vspace{8pt}
\clearpage

\section{归园田居}

\beisong{泛读}

\begin{center}
    {\zihao{-0}\bfseries 归园田居}

    \vspace{18pt}

    \begin{large}\xuwen
        〔陶渊明〕
    \end{large}

\end{center}

    \vspace{18pt}  % 起首

\begin{center}
    \begin{Large}

        少无适俗韵，性本爱丘山。

        误落尘网中，一去三十年。

        羁鸟恋旧林，池鱼思故渊。

        开荒南野际，守拙归园田。

        方宅十余亩，草屋八九间。

        榆柳荫后檐，桃李罗堂前。

        暧暧远人村，依依墟里烟。

        狗吠深巷中，鸡鸣桑树颠。

        户庭无尘杂，虚室有余闲。

        久在樊笼里，复得返自然。

    \end{Large}

\end{center}

\vspace{8pt}
\clearpage

\section{拟行路难}

\beisong{泛读}

\begin{center}
    {\zihao{-0}\bfseries 拟行路难}

    \vspace{18pt}

    \begin{large}\xuwen
        〔鲍照〕
    \end{large}

\end{center}

    \vspace{18pt}  % 起首

\begin{center}
    \begin{Large}

        泻水置平地，各自东西南北流。

        人生亦有命，安能行叹复坐愁？

        酌酒以自宽，举杯断绝歌路难。

        心非木石岂无感？吞声踯躅不敢言。\hspace{24pt}

    \end{Large}

\end{center}

\vspace{8pt}
\clearpage

\section{书愤}

\beisong{背诵}

\begin{center}
    {\zihao{-0}\bfseries 书愤}

    \vspace{18pt}

    \begin{large}\xuwen
        〔陆游〕
    \end{large}

\end{center}

    \vspace{18pt}  % 起首

\begin{center}
    \begin{Large}

        早岁那知世事艰，中原北望气如山。

        楼船夜雪瓜洲渡，铁马秋风大散关。

        塞上长城空自许，镜中衰鬓已先斑。

        出师一表真名世，千载谁堪伯仲间。

    \end{Large}

\end{center}

\vspace{8pt}
\clearpage

\section{永遇乐·京口北固亭怀古}

\beisong{背诵}

\begin{center}
    {\zihao{-0}\bfseries 永遇乐·京口北固亭怀古}

    \vspace{18pt}

    \begin{large}\xuwen
        〔辛弃疾〕
    \end{large}

\end{center}

    \vspace{18pt}  % 起首

\begin{center}
    \begin{Large}

        千古江山，英雄无觅孙仲谋处。\hspace{42pt}

        舞榭歌台，风流总被雨打风吹去。\hspace{24pt}

        斜阳草树，寻常巷陌，人道寄奴曾住。

        \hspace{18pt}想当年，金戈铁马，气吞万里如虎。

        \vspace{18pt}  % 分阙

        元嘉草草，封狼居胥，赢得仓皇北顾。

        四十三年，望中犹记，烽火扬州路。\hspace{5pt}

        可堪回首，佛狸祠下，一片神鸦社鼓。

        凭谁问：廉颇老矣，尚能饭否？\hspace{5pt}

    \end{Large}

\end{center}

\vspace{8pt}
\clearpage

\section{闻乐天授江州司马}

\beisong{背诵}

\begin{center}
    {\zihao{-0}\bfseries 闻乐天授江州司马}

    \vspace{18pt}

    \begin{large}\xuwen
        〔元稹〕
    \end{large}

\end{center}

    \vspace{18pt}  % 起首

\begin{center}
    \begin{Large}

        残灯无焰影幢幢，此夕闻君谪九江。

        垂死病中惊坐起，暗风吹雨入寒窗。

    \end{Large}

\end{center}

\vspace{8pt}
\clearpage

\section{逢入京使}

\beisong{泛读}

\begin{center}
    {\zihao{-0}\bfseries 逢入京使}

    \vspace{18pt}

    \begin{large}\xuwen
        〔岑参〕
    \end{large}

\end{center}

    \vspace{18pt}

\begin{center}
    \begin{Large}

        故园东望路漫漫，双袖龙钟泪不干。

        马上相逢无纸笔，凭君传语报平安。

    \end{Large}

\end{center}

\vspace{8pt}
\clearpage

\section{新安吏}

\beisong{泛读}

\begin{center}
    {\zihao{-0}\bfseries 新安吏}

    \vspace{18pt}

    \begin{large}\xuwen
        〔杜甫〕
    \end{large}

\end{center}

    \vspace{18pt}  % 起首

\begin{center}
    \begin{Large}

        客行新安道，喧呼闻点兵。

        借问新安吏，县小更无丁。

        府帖昨夜下，次选中男行。

        中男绝短小，何以守王城？

        肥男有母送，瘦男独伶俜。

        白水暮东流，青山犹哭声。

        莫自使眼枯，收汝泪纵横。

        眼枯即见骨，天地终无情。

        我军取相州，日夕望其平。

        岂意贼难料，归军星散营。

        就粮近故垒，练卒依旧京。

        掘壕不到水，牧马役亦轻。

        况乃王师顺，抚养甚分明。

        送行勿泣血，仆射如父兄。

    \end{Large}

\end{center}

\vspace{8pt}
\clearpage

\chapter{第九级}
\clearpage

\section{月下独酌}

\beisong{背诵}

\begin{center}
    {\zihao{-0}\bfseries 月下独酌}

    \vspace{18pt}

    \begin{large}\xuwen
        〔李白〕
    \end{large}

\end{center}

    \vspace{18pt}

\begin{center}
    \begin{Large}

        花间一壶酒，独酌无相亲。

        举杯邀明月，对影成三人。

        月既不解饮，影徒随我身。

        暂伴月将影，行乐须及春。

        我歌月徘徊，我舞影零乱。

        醒时同交欢，醉后各分散。

        永结无情游，相期邈云汉。

    \end{Large}

\end{center}

\vspace{8pt}
\clearpage

\section{行宫}

\beisong{泛读}

\begin{center}
    {\zihao{-0}\bfseries 行宫}

    \vspace{18pt}

    \begin{large}\xuwen
        〔元稹〕
    \end{large}

\end{center}

    \vspace{18pt}

\begin{center}
    \begin{Large}

        寥落古行宫，宫花寂寞红。

        白头宫女在，闲坐说玄宗。

    \end{Large}

\end{center}

\vspace{8pt}
\clearpage

\section{春江花月夜}

\beisong{精读}

\begin{center}
    {\zihao{-0}\bfseries 春江花月夜}

    \vspace{18pt}

    \begin{large}\xuwen
        〔张若虚〕
    \end{large}

\end{center}

    \vspace{10pt}  % 起首

\begin{center}
    \begin{Large}

        春江潮水连海平，海上明月共潮生。

        滟滟随波千万里，何处春江无月明。

        江流宛转绕芳甸，月照花林皆似霰；

        空里流霜不觉飞，汀上白沙看不见。

        江天一色无纤尘，皎皎空中孤月轮。

        江畔何人初见月？江月何年初照人？

        人生代代无穷已，江月年年只相似。

        不知江月待何人，但见长江送流水。

        白云一片去悠悠，青枫浦上不胜愁。

        谁家今夜扁舟子？何处相思明月楼？

        可怜楼上月徘徊，应照离人妆镜台。

        玉户帘中卷不去，捣衣砧上拂还来。

        此时相望不相闻，愿逐月华流照君。

        鸿雁长飞光不度，鱼龙潜跃水成文。

        昨夜闲潭梦落花，可怜春半不还家。

        江水流春去欲尽，江潭落月复西斜。

        斜月沉沉藏海雾，碣石潇湘无限路。

        不知乘月几人归，落月摇情满江树。

    \end{Large}

\end{center}

\vspace{8pt}
\clearpage

\section{蜀道难}

\beisong{背诵}

\begin{center}
    {\zihao{-0}\bfseries 蜀道难}

    \vspace{18pt}

    \begin{large}\xuwen
        〔李白〕
    \end{large}

\end{center}

    \vspace{18pt}  % 起首

\begin{center}
    \begin{Large}

        \hspace{16pt}噫吁嚱，危乎高哉！
        
        \hspace{16pt}蜀道之难，难于上青天！

        蚕丛及鱼凫，开国何茫然！

        尔来四万八千岁，不与秦塞通人烟。

        西当太白有鸟道，可以横绝峨眉巅。

        \hspace{34pt}地崩山摧壮士死，然后天梯石栈相钩连。

        上有六龙回日之高标，下有冲波逆折之回川。

        黄鹤之飞尚不得过，猿猱欲度愁攀援。

        \hspace{34pt}青泥何盘盘，百步九折萦岩峦。

        扪参历井仰胁息，以手抚膺坐长叹。

        问君西游何时还？畏途巉岩不可攀。

        但见悲鸟号古木，雄飞雌从绕林间。

        又闻子规啼夜月，愁空山。\hspace{58pt}

        \hspace{16pt}蜀道之难，难于上青天！
        
        使人听此凋朱颜！\phantom{日日日日日日日}\hspace{6pt}

        连峰去天不盈尺，枯松倒挂倚绝壁。

        飞湍瀑流争喧豗，砯崖转石万壑雷。

        其险也如此，嗟尔远道之人胡为乎来哉！

        剑阁峥嵘而崔嵬，一夫当关，万夫莫开。

        所守或匪亲，化为狼与豺。

        朝避猛虎，夕避长蛇；
        
        磨牙吮血，杀人如麻。

        锦城虽云乐，不如早还家。

        \hspace{16pt}蜀道之难，难于上青天！
        
        侧身西望长咨嗟！\phantom{日日日日日日日}\hspace{6pt}

    \end{Large}

\end{center}

\vspace{8pt}
\clearpage

\section{兵车行}

\beisong{精读}

\begin{center}
    {\zihao{-0}\bfseries 兵车行}

    \vspace{18pt}

    \begin{large}\xuwen
        〔杜甫〕
    \end{large}

\end{center}

    \vspace{18pt}  % 起首

\begin{center}
    \begin{Large}

        车辚辚，马萧萧，行人弓箭各在腰。

        爷娘妻子走相送，尘埃不见咸阳桥。

        牵衣顿足拦道哭，哭声直上干云霄。

        道旁过者问行人，行人但云点行频。

        或从十五北防河，便至四十西营田。

        去时里正与裹头，归来头白还戍边。

        边庭流血成海水，武皇开边意未已。

        君不闻汉家山东二百州，千村万落生荆杞。\hspace{40pt}

        纵有健妇把锄犁，禾生陇亩无东西。

        况复秦兵耐苦战，被驱不异犬与鸡。

        长者虽有问，役夫敢申恨？

        且如今年冬，未休关西卒。

        县官急索租，租税从何出？

        信知生男恶，反是生女好。

        生女犹得嫁比邻，生男埋没随百草。

        君不见，青海头，古来白骨无人收。

        新鬼烦冤旧鬼哭，天阴雨湿声啾啾！

    \end{Large}

\end{center}

\vspace{8pt}
\clearpage

\section{燕歌行}

\beisong{精读}

\begin{center}
    {\zihao{-0}\bfseries 燕歌行}

    \vspace{18pt}

    \begin{large}\xuwen
        〔高适〕
    \end{large}

\end{center}

    \vspace{18pt}  % 起首

\begin{center}
    \begin{Large}

        汉家烟尘在东北， 汉将辞家破残贼。

        男儿本自重横行， 天子非常赐颜色。

        摐金伐鼓下榆关， 旌旆逶迤碣石间。

        校尉羽书飞瀚海， 单于猎火照狼山。

        山川萧条极边土， 胡骑凭陵杂风雨。

        战士军前半死生， 美人帐下犹歌舞！

        大漠穷秋塞草腓， 孤城落日斗兵稀。

        身当恩遇恒轻敌， 力尽关山未解围。

        铁衣远戍辛勤久， 玉箸应啼别离后。

        少妇城南欲断肠， 征人蓟北空回首。

        边庭飘飖那可度， 绝域苍茫更何有！

        杀气三时作阵云， 寒声一夜传刁斗。

        相看白刃血纷纷， 死节从来岂顾勋？

        君不见沙场征战苦，至今犹忆李将军！

    \end{Large}

\end{center}

\vspace{8pt}
\clearpage

\section{客至}

\beisong{背诵}

\begin{center}
    {\zihao{-0}\bfseries 客至}

    \vspace{18pt}

    \begin{large}\xuwen
        〔杜甫〕
    \end{large}

\end{center}

    \vspace{18pt}  % 起首

\begin{center}
    \begin{Large}

        舍南舍北皆春水，但见群鸥日日来。

        花径不曾缘客扫，蓬门今始为君开。

        盘餐市远无兼味，樽酒家贫只旧醅。

        肯与邻翁相对饮，隔篱呼取尽馀杯。

    \end{Large}

\end{center}

\vspace{8pt}
\clearpage

\section{虞美人}

\beisong{背诵}

\begin{center}
    {\zihao{-0}\bfseries 虞美人}

    \vspace{18pt}

    \begin{large}\xuwen
        〔李煜〕
    \end{large}

\end{center}

    \vspace{18pt}  % 起首

\begin{center}
    \begin{Large}

        春花秋月何时了？往事知多少。\hspace{58pt}

        小楼昨夜又东风，故国不堪回首月明中。

        \vspace{18pt}  % 分阙

        雕栏玉砌应犹在，只是朱颜改。\hspace{58pt}

        问君能有几多愁？恰似一江春水向东流。

    \end{Large}

\end{center}

\vspace{8pt}
\clearpage

\section{望海潮}

\beisong{精读}

\begin{center}
    {\zihao{-0}\bfseries 望海潮}

    \vspace{18pt}

    \begin{large}\xuwen
        〔柳永〕
    \end{large}

\end{center}

    \vspace{18pt}  % 起首

\begin{center}
    \begin{Large}

        东南形胜，三吴都会，

        钱塘自古繁华，

        烟柳画桥，风帘翠幕，

        参差十万人家。

        云树绕堤沙，怒涛卷霜雪，天堑无涯。

        市列珠玑，户盈罗绮，竞豪奢。

        \vspace{18pt}  % 分阙

        重湖叠巘清嘉。

        有三秋桂子，十里荷花。

        羌管弄晴，菱歌泛夜，

        嬉嬉钓叟莲娃。

        千骑拥高牙。乘醉听箫鼓，吟赏烟霞。

        异日图将好景，归去凤池夸。

    \end{Large}

\end{center}

\vspace{8pt}
\clearpage

\section{遣怀}

\beisong{泛读}

\begin{center}
    {\zihao{-0}\bfseries 遣怀}

    \vspace{18pt}

    \begin{large}\xuwen
        〔杜牧〕
    \end{large}

\end{center}

    \vspace{18pt}

\begin{center}
    \begin{Large}

        落魄江湖载酒行，楚腰纤细掌中轻。

        十年一觉扬州梦，赢得青楼薄幸名。

    \end{Large}

\end{center}

\vspace{8pt}
\clearpage

\section{声声慢·寻寻觅觅}

\beisong{背诵}

\begin{center}
    {\zihao{-0}\bfseries 声声慢·寻寻觅觅}

    \vspace{18pt}

    \begin{large}\xuwen
        〔李清照〕
    \end{large}

\end{center}

    \vspace{18pt}  % 起首

\begin{center}
    \begin{Large}

        寻寻觅觅，冷冷清清，

        凄凄惨惨戚戚。

        乍暖还寒时候，最难将息。

        三杯两盏淡酒，怎敌他、晚来风急？

        雁过也，正伤心，却是旧时相识。

        \vspace{18pt}  % 分阙

        满地黄花堆积，憔悴损，

        如今有谁堪摘？

        守着窗儿，独自怎生得黑？

        梧桐更兼细雨，到黄昏、点点滴滴。

        这次第，怎一个愁字了得！

    \end{Large}

\end{center}

\vspace{8pt}
\clearpage

\section{阁夜}

\beisong{泛读}

\begin{center}
    {\zihao{-0}\bfseries 阁夜}

    \vspace{18pt}

    \begin{large}\xuwen
        〔杜甫〕
    \end{large}

\end{center}

    \vspace{18pt}  % 起首

\begin{center}
    \begin{Large}

        岁暮阴阳催短景，天涯霜雪霁寒宵。

        五更鼓角声悲壮，三峡星河影动摇。

        野哭千家闻战伐，夷歌数处起渔樵。

        卧龙跃马终黄土，人事音书漫寂寥。

    \end{Large}

\end{center}

\vspace{8pt}
\clearpage

\section{潼关吏}

\beisong{泛读}

\begin{center}
    {\zihao{-0}\bfseries 潼关吏}

    \vspace{18pt}

    \begin{large}\xuwen
        〔杜甫〕
    \end{large}

\end{center}

    \vspace{18pt}  % 起首

\begin{center}
    \begin{Large}

        士卒何草草，筑城潼关道。

        大城铁不如，小城万丈馀。

        借问潼关吏：修关还备胡？

        要我下马行，为我指山隅：

        连云列战格，飞鸟不能逾。

        胡来但自守，岂复忧西都。

        丈人视要处，窄狭容单车。

        艰难奋长戟，万古用一夫。

        哀哉桃林战，百万化为鱼。

        请嘱防关将，慎勿学哥舒！

    \end{Large}

\end{center}

\vspace{8pt}
\clearpage

\section{诗经·硕鼠}

\beisong{泛读}

\begin{center}
    {\zihao{-0}\bfseries 诗经·硕鼠}

    \vspace{18pt}

    \begin{large}\xuwen
        〔诗经〕
    \end{large}

\end{center}

    \vspace{18pt}  % 起首

\begin{center}
    \begin{Large}

        硕鼠硕鼠，无食我黍！

        三岁贯女，莫我肯顾。

        逝将去女，适彼乐土。

        乐土乐土，爰得我所。

        \vspace{18pt}  % 分阙

        硕鼠硕鼠，无食我麦！

        三岁贯女，莫我肯德。

        逝将去女，适彼乐国。

        乐国乐国，爰得我直。

        \vspace{18pt}  % 分阙

        硕鼠硕鼠，无食我苗！

        三岁贯女，莫我肯劳。

        逝将去女，适彼乐郊。

        乐郊乐郊，谁之永号？

    \end{Large}

\end{center}

\vspace{8pt}
\clearpage

\section{丹青引赠曹将军霸}

\beisong{泛读}

\begin{center}
    {\zihao{-0}\bfseries 丹青引赠曹将军霸}

    \vspace{18pt}

    \begin{large}\xuwen
        〔杜甫〕
    \end{large}

\end{center}

    \vspace{18pt}  % 起首

\begin{center}
    \begin{Large}

        将军魏武之子孙，于今为庶为清门。

        英雄割据虽已矣，文采风流今尚存。

        学书初学卫夫人，但恨无过王右军。

        丹青不知老将至，富贵于我如浮云。

        开元之中常引见，承恩数上南薰殿。

        凌烟功臣少颜色，将军下笔开生面。

        良相头上进贤冠，猛将腰间大羽箭。

        褒公鄂公毛发动，英姿飒爽来酣战。

        先帝天马玉花骢，画工如山貌不同。

        是日牵来赤墀下，迥立阊阖生长风。

        诏谓将军拂绢素，意匠惨澹经营中。

        斯须九重真龙出，一洗万古凡马空。

        玉花却在御榻上，榻上庭前屹相向。

        至尊含笑催赐金，圉人太仆皆惆怅。

        弟子韩干早入室，亦能画马穷殊相。

        干惟画肉不画骨，忍使骅骝气凋丧。

        将军画善盖有神，必逢佳士亦写真。

        即今漂泊干戈际，屡貌寻常行路人。

        途穷反遭俗眼白，世上未有如公贫。

        但看古来盛名下，终日坎壈缠其身。

    \end{Large}

\end{center}

\vspace{8pt}
\clearpage

\chapter{第十级}
\clearpage

\section{长恨歌}

\beisong{泛读}

\begin{center}
    {\zihao{-0}\bfseries 长恨歌}

    \vspace{18pt}

    \begin{large}\xuwen
        〔白居易〕
    \end{large}

\end{center}

    \vspace{18pt}  % 起首

\begin{center}
    \begin{Large}

        汉皇重色思倾国，御宇多年求不得。

        杨家有女初长成，养在深闺人未识。

        天生丽质难自弃，一朝选在君王侧。

        回眸一笑百媚生，六宫粉黛无颜色。

        春寒赐浴华清池，温泉水滑洗凝脂。

        侍儿扶起娇无力，始是新承恩泽时。

        云鬓花颜金步摇，芙蓉帐暖度春宵。

        春宵苦短日高起，从此君王不早朝。

        承欢侍宴无闲暇，春从春游夜专夜。

        后宫佳丽三千人，三千宠爱在一身。

        金屋妆成娇侍夜，玉楼宴罢醉和春。

        姊妹弟兄皆列土，可怜光彩生门户。

        遂令天下父母心，不重生男重生女。

        骊宫高处入青云，仙乐风飘处处闻。

        缓歌慢舞凝丝竹，尽日君王看不足。

        渔阳鼙鼓动地来，惊破霓裳羽衣曲。

        九重城阙烟尘生，千乘万骑西南行。

        翠华摇摇行复止，西出都门百余里。

        六军不发无奈何，宛转蛾眉马前死。

        花钿委地无人收，翠翘金雀玉搔头。

        君王掩面救不得，回看血泪相和流。

        黄埃散漫风萧索，云栈萦纡登剑阁。

        峨嵋山下少人行，旌旗无光日色薄。

        蜀江水碧蜀山青，圣主朝朝暮暮情。

        行宫见月伤心色，夜雨闻铃肠断声。

        天旋地转回龙驭，到此踌躇不能去。

        马嵬坡下泥土中，不见玉颜空死处。

        君臣相顾尽沾衣，东望都门信马归。

        归来池苑皆依旧，太液芙蓉未央柳。

        芙蓉如面柳如眉，对此如何不泪垂？

        春风桃李花开日，秋雨梧桐叶落时。

        西宫南内多秋草，落叶满阶红不扫。

        梨园弟子白发新，椒房阿监青娥老。

        夕殿萤飞思悄然，孤灯挑尽未成眠。

        迟迟钟鼓初长夜，耿耿星河欲曙天。

        鸳鸯瓦冷霜华重，翡翠衾寒谁与共？

        悠悠生死别经年，魂魄不曾来入梦。

        临邛道士鸿都客，能以精诚致魂魄。

        为感君王辗转思，遂教方士殷勤觅。

        排空驭气奔如电，升天入地求之遍。

        上穷碧落下黄泉，两处茫茫皆不见。

        忽闻海上有仙山，山在虚无缥缈间。

        楼阁玲珑五云起，其中绰约多仙子。

        中有一人字太真，雪肤花貌参差是。

        金阙西厢叩玉扃，转教小玉报双成。

        闻道汉家天子使，九华帐里梦魂惊。

        揽衣推枕起徘徊，珠箔银屏迤逦开。

        云鬓半偏新睡觉，花冠不整下堂来。

        风吹仙袂飘飖举，犹似霓裳羽衣舞。

        玉容寂寞泪阑干，梨花一枝春带雨。

        含情凝睇谢君王，一别音容两渺茫。

        昭阳殿里恩爱绝，蓬莱宫中日月长。

        回头下望人寰处，不见长安见尘雾。

        惟将旧物表深情，钿合金钗寄将去。

        钗留一股合一扇，钗擘黄金合分钿。

        但教心似金钿坚，天上人间会相见。

        临别殷勤重寄词，词中有誓两心知。

        七月七日长生殿，夜半无人私语时。

        在天愿作比翼鸟，在地愿为连理枝。

        天长地久有时尽，此恨绵绵无绝期。

    \end{Large}

\end{center}

\vspace{8pt}
\clearpage

\section{庐山谣寄卢侍御虚舟}

\beisong{泛读}

\begin{center}
    {\zihao{-0}\bfseries 庐山谣寄卢侍御虚舟}

    \vspace{18pt}

    \begin{large}\xuwen
        〔李白〕
    \end{large}

\end{center}

    \vspace{18pt}

\begin{center}
    \begin{Large}

        我本楚狂人，凤歌笑孔丘。

        手持绿玉杖，朝别黄鹤楼。

        五岳寻仙不辞远，一生好入名山游。

        庐山秀出南斗傍，屏风九叠云锦张。

        影落明湖青黛光，金阙前开二峰长。

        银河倒挂三石梁，香炉瀑布遥相望。

        回崖沓障凌苍苍。

        翠影红霞映朝日，鸟飞不到吴天长。

        登高壮观天地间，大江茫茫去不黄。

        黄云万里动风色，白波九道流雪山。

        好为庐山谣，兴因庐山发。

        闲窥石镜清我心，谢公行处苍苔没。

        早服还丹无世情，琴心三叠道初成。

        遥见仙人彩云里，手把芙蓉朝玉京。

        先期汗漫九垓上，愿接卢敖游太清。

    \end{Large}

\end{center}

\vspace{8pt}
\clearpage

\section{梦游天姥吟留别}

\beisong{背诵}

\begin{center}
    {\zihao{-0}\bfseries 梦游天姥吟留别}

    \vspace{18pt}

    \begin{large}\xuwen
        〔李白〕
    \end{large}

\end{center}

    \vspace{18pt}  % 起首

\begin{center}
    \begin{Large}

        \hspace{34pt}海客谈瀛洲，烟涛微茫信难求。

        \hspace{34pt}越人语天姥，云霞明灭或可睹。

        天姥连天向天横，势拔五岳掩赤城。

        天台四万八千丈，对此欲倒东南倾。

        我欲因之梦吴越，一夜飞度镜湖月。

        湖月照我影，送我至剡溪。

        谢公宿处今尚在，渌水荡漾清猿啼。

        脚著谢公屐，身登青云梯。

        半壁见海日，空中闻天鸡。

        千岩万转路不定，迷花倚石忽已暝。

        熊咆龙吟殷岩泉，栗深林兮惊层巅。

        云青青兮欲雨，水澹澹兮生烟。

        列缺霹雳，丘峦崩摧。

        洞天石扉，訇然中开。

        青冥浩荡不见底，日月照耀金银台。

        \hspace{34pt}霓为衣兮风为马，云之君兮纷纷而来下。

        虎鼓瑟兮鸾回车，仙之人兮列如麻。

        忽魂悸以魄动，恍惊起而长嗟。

        惟觉时之枕席，失向来之烟霞。

        世间行乐亦如此，古来万事东流水。

        别君去兮何时还？

        且放白鹿青崖间。须行即骑访名山。

        安能摧眉折腰事权贵，使我不得开心颜！\hspace{24pt}

    \end{Large}

\end{center}

\vspace{8pt}
\clearpage

\section{将进酒}

\beisong{背诵}

\begin{center}
    {\zihao{-0}\bfseries 将进酒}

    \vspace{18pt}

    \begin{large}\xuwen
        〔李白〕
    \end{large}

\end{center}

    \vspace{18pt}  % 起首

\begin{center}
    \begin{Large}

        君不见黄河之水天上来，奔流到海不复回。\hspace{42pt}

        君不见高堂明镜悲白发，朝如青丝暮成雪。\hspace{42pt}

        人生得意须尽欢，莫使金樽空对月。

        天生我材必有用，千金散尽还复来。

        烹羊宰牛且为乐，会须一饮三百杯。

        岑夫子，丹丘生，将进酒，杯莫停。

        \hspace{36pt}与君歌一曲，请君为我倾耳听。

        钟鼓馔玉不足贵，但愿长醉不复醒。

        古来圣贤皆寂寞，惟有饮者留其名。

        陈王昔时宴平乐，斗酒十千恣欢谑。

        主人何为言少钱，径须沽取对君酌。

        五花马，千金裘，

        呼儿将出换美酒，与尔同销万古愁。

    \end{Large}

\end{center}

\vspace{8pt}
\clearpage

\section{题金陵渡}

\beisong{泛读}

\begin{center}
    {\zihao{-0}\bfseries 题金陵渡}

    \vspace{18pt}

    \begin{large}\xuwen
        〔张祜〕
    \end{large}

\end{center}

    \vspace{18pt}

\begin{center}
    \begin{Large}

        金陵津渡小山楼，一宿行人自可愁。

        潮落夜江斜月里，两三星火是瓜州。

    \end{Large}

\end{center}

\vspace{8pt}
\clearpage

\section{蜀相}

\beisong{背诵}

\begin{center}
    {\zihao{-0}\bfseries 蜀相}

    \vspace{18pt}

    \begin{large}\xuwen
        〔杜甫〕
    \end{large}

\end{center}

    \vspace{18pt}  % 起首

\begin{center}
    \begin{Large}

        丞相祠堂何处寻，锦官城外柏森森。

        映阶碧草自春色，隔叶黄鹂空好音。

        三顾频烦天下计，两朝开济老臣心。

        出师未捷身先死，长使英雄泪满襟。

    \end{Large}

\end{center}

\vspace{8pt}
\clearpage

\section{登高}

\beisong{背诵}

\begin{center}
    {\zihao{-0}\bfseries 登高}

    \vspace{18pt}

    \begin{large}\xuwen
        〔杜甫〕
    \end{large}

\end{center}

    \vspace{12pt}\mbox{}

    \begin{large}\xuwen
        开元二十六年，客有从元戎出塞而还者，作《燕歌行》以示，适感征戍之事，因而和焉。
    \end{large}


    \vspace{8pt}

\begin{center}
    \begin{Large}

        风急天高猿啸哀，渚清沙白鸟飞回。

        无边落木萧萧下，不尽长江滚滚来。

        万里悲秋常作客，百年多病独登台。

        艰难苦恨繁霜鬓，潦倒新停浊酒杯。

    \end{Large}

\end{center}

\vspace{8pt}
\clearpage

\section{送崔九}

\beisong{泛读}

\begin{center}
    {\zihao{-0}\bfseries 送崔九}

    \vspace{18pt}

    \begin{large}\xuwen
        〔裴迪〕
    \end{large}

\end{center}

    \vspace{18pt}

\begin{center}
    \begin{Large}

        归山深浅去，须尽丘壑美。

        莫学武陵人，暂游桃源里。

    \end{Large}

\end{center}

\vspace{8pt}
\clearpage

\section{遣悲怀}

\beisong{泛读}

\begin{center}
    {\zihao{-0}\bfseries 遣悲怀}

    \vspace{18pt}

    \begin{large}\xuwen
        〔元稹〕
    \end{large}

\end{center}

    \vspace{18pt}

\begin{center}
    \begin{Large}

        昔日戏言身后事，今朝都到眼前来。

        衣裳已施行看尽，针线犹存未忍开。

        尚想旧情怜婢仆，也曾因梦送钱财。

        诚知此恨人人有，贫贱夫妻百事哀。

    \end{Large}

\end{center}

\vspace{8pt}
\clearpage

\section{鹊桥仙·纤云弄巧}

\beisong{背诵}

\begin{center}
    {\zihao{-0}\bfseries 鹊桥仙·纤云弄巧}

    \vspace{18pt}

    \begin{large}\xuwen
        〔秦观〕
    \end{large}

\end{center}

    \vspace{18pt}  % 起首

\begin{center}
    \begin{Large}

        纤云弄巧，飞星传恨，银汉迢迢暗度。

        金风玉露一相逢，便胜却人间无数。

        \vspace{18pt}  % 分阙

        柔情似水，佳期如梦，忍顾鹊桥归路。

        两情若是久长时，又岂在朝朝暮暮。

    \end{Large}

\end{center}

\vspace{8pt}
\clearpage

\section{扬州慢·淮左名都}

\beisong{背诵}

\begin{center}
    {\zihao{-0}\bfseries 扬州慢·淮左名都}

    \vspace{18pt}

    \begin{large}\xuwen
        〔姜夔〕
    \end{large}

\end{center}

    \vspace{18pt}  % 起首

\begin{center}
    \begin{Large}

        淮左名都，竹西佳处，解鞍少驻初程。

        过春风十里。尽荠麦青青。

        自胡马窥江去后，

        废池乔木，犹厌言兵。

        渐黄昏，清角吹寒。都在空城。

        \vspace{18pt}  % 分阙

        杜郎俊赏，算而今、重到须惊。

        纵豆蔻词工，青楼梦好，难赋深情。

        二十四桥仍在，波心荡、冷月无声。

        念桥边红药，年年知为谁生。

    \end{Large}

\end{center}

\vspace{8pt}
\clearpage

\section{菩萨蛮·人人尽说江南好}

\beisong{泛读}

\begin{center}
    {\zihao{-0}\bfseries 菩萨蛮·人人尽说江南好}

    \vspace{18pt}

    \begin{large}\xuwen
        〔韦庄〕
    \end{large}

\end{center}

    \vspace{18pt}  % 起首

\begin{center}
    \begin{Large}

        人人尽说江南好，

        游人只合江南老。

        春水碧于天，画船听雨眠。

        \vspace{18pt}  % 分阙

        垆边人似月，

        皓腕凝霜雪。

        未老莫还乡，还乡须断肠。

    \end{Large}

\end{center}

\vspace{8pt}
\clearpage

\section{新婚别}

\beisong{泛读}

\begin{center}
    {\zihao{-0}\bfseries 新婚别}

    \vspace{18pt}

    \begin{large}\xuwen
        〔杜甫〕
    \end{large}

\end{center}

    \vspace{18pt}  % 起首

\begin{center}
    \begin{Large}

        兔丝附蓬麻，引蔓故不长。

        嫁女与征夫，不如弃路旁。

        结发为妻子，席不暖君床。

        暮婚晨告别，无奈太匆忙!

        君行虽不远，守边赴河阳。

        妾身未分明，何以拜姑嫜？

        父母养我时，日夜令我藏。

        生女有所归，鸡狗亦得将。

        君今往死地，沉痛迫中肠。

        誓欲随君去，形势反苍黄。

        勿为新婚念，努力事戎行。

        妇人在军中，兵气恐不扬。

        自嗟贫家女，久致罗襦裳。

        罗襦不复施，对君洗红妆。

        仰视百鸟飞，大小必双翔。

        人事多错迕，与君永相望！

    \end{Large}

\end{center}

\vspace{8pt}
\clearpage

\section{诗经·伐檀}

\beisong{泛读}

\begin{center}
    {\zihao{-0}\bfseries 诗经·伐檀}

    \vspace{18pt}

    \begin{large}\xuwen
        〔诗经〕
    \end{large}

\end{center}

    \vspace{18pt}  % 起首

\begin{center}
    \begin{Large}

        坎坎伐檀兮，置之河之干兮。

        河水清且涟猗。

        不稼不穑，胡取禾三百廛兮？

        不狩不猎，胡瞻尔庭有县貆兮？

        彼君子兮，不素餐兮！

        \vspace{18pt}  % 分阙

        坎坎伐辐兮，置之河之侧兮。

        河水清且直猗。

        不稼不穑，胡取禾三百繶兮？

        不狩不猎，胡瞻尔庭有县特兮？

        彼君子兮，不素食兮！

        \vspace{18pt}  % 分阙

        坎坎伐轮兮，置之河之漘兮。

        河水清且沦猗。

        不稼不穑，胡取禾三百囷兮？

        不狩不猎，胡瞻尔庭有县鹑兮？

        彼君子兮，不素飧兮！

    \end{Large}

\end{center}

\vspace{8pt}
\clearpage

\chapter{第十一级}
\clearpage

\section{李凭箜篌引}

\beisong{泛读}

\begin{center}
    {\zihao{-0}\bfseries 李凭箜篌引}

    \vspace{18pt}

    \begin{large}\xuwen
        〔李贺〕
    \end{large}

\end{center}

    \vspace{18pt}  % 起首

\begin{center}
    \begin{Large}

        吴丝蜀桐张高秋，空山凝云颓不流。

        江娥啼竹素女愁，李凭中国弹箜篌。

        昆山玉碎凤凰叫，芙蓉泣露香兰笑。

        十二门前融冷光，二十三丝动紫皇。

        女娲炼石补天处，石破天惊逗秋雨。

        梦入神山教神妪，老鱼跳波瘦蛟舞。

        吴质不眠倚桂树，露脚斜飞湿寒兔。

    \end{Large}

\end{center}

\vspace{8pt}
\clearpage

\section{菩萨蛮}

\beisong{泛读}

\begin{center}
    {\zihao{-0}\bfseries 菩萨蛮}

    \vspace{18pt}

    \begin{large}\xuwen
        〔温庭钧〕
    \end{large}

\end{center}

    \vspace{18pt}  % 起首

\begin{center}
    \begin{Large}

        小山重叠金明灭，鬓云欲度香腮雪。

        懒起画蛾眉，弄妆梳洗迟。

        照花前后镜，花面交相映。

        新帖绣罗襦，双双金鹧鸪。

    \end{Large}

\end{center}

\vspace{8pt}
\clearpage

\section{锦瑟}

\beisong{背诵}

\begin{center}
    {\zihao{-0}\bfseries 锦瑟}

    \vspace{18pt}

    \begin{large}\xuwen
        〔李商隐〕
    \end{large}

\end{center}

    \vspace{18pt}  % 起首

\begin{center}
    \begin{Large}

        锦瑟无端五十弦，一弦一柱思华年。

        庄生晓梦迷蝴蝶，望帝春心托杜鹃。

        沧海月明珠有泪，蓝田日暖玉生烟。

        此情可待成追忆，只是当时已惘然。

    \end{Large}

\end{center}

\vspace{8pt}
\clearpage

\section{夜上受降城闻笛}

\beisong{泛读}

\begin{center}
    {\zihao{-0}\bfseries 夜上受降城闻笛}

    \vspace{18pt}

    \begin{large}\xuwen
        〔李益〕
    \end{large}

\end{center}

    \vspace{18pt}

\begin{center}
    \begin{Large}

        回乐峰前沙似雪，受降城外月如霜。

        不知何处吹芦管，一夜征人尽望乡。

    \end{Large}

\end{center}

\vspace{8pt}
\clearpage

\section{登楼}

\beisong{泛读}

\begin{center}
    {\zihao{-0}\bfseries 登楼}

    \vspace{18pt}

    \begin{large}\xuwen
        〔杜甫〕
    \end{large}

\end{center}

    \vspace{18pt}

\begin{center}
    \begin{Large}

        花近高楼伤客心，万方多难此登临。

        锦江春色来天地，玉垒浮云变古今。

        北极朝庭终不改，西山寇盗莫相侵！

        可怜后主还祠庙，日暮聊为梁父吟。

    \end{Large}

\end{center}

\vspace{8pt}
\clearpage

\section{征人怨}

\beisong{泛读}

\begin{center}
    {\zihao{-0}\bfseries 征人怨}

    \vspace{18pt}

    \begin{large}\xuwen
        〔柳中庸〕
    \end{large}

\end{center}

    \vspace{18pt}

\begin{center}
    \begin{Large}

        岁岁金河复玉关，朝朝马策与刀环。

        三春白雪归青冢，万里黄河绕黑山。

    \end{Large}

\end{center}

\vspace{8pt}
\clearpage

\section{闺怨}

\beisong{泛读}

\begin{center}
    {\zihao{-0}\bfseries 闺怨}

    \vspace{18pt}

    \begin{large}\xuwen
        〔王昌龄〕
    \end{large}

\end{center}

    \vspace{18pt}

\begin{center}
    \begin{Large}

        闺中少妇不知愁，春日凝妆上翠楼。

        忽见陌头杨柳色，悔教夫婿觅封侯。

    \end{Large}

\end{center}

\vspace{8pt}
\clearpage

\section{念奴娇·赤壁怀古}

\beisong{背诵}

\begin{center}
    {\zihao{-0}\bfseries 念奴娇·赤壁怀古}

    \vspace{18pt}

    \begin{large}\xuwen
        〔苏轼〕
    \end{large}

\end{center}

    \vspace{18pt}  % 起首

\begin{center}
    \begin{Large}

        大江东去，浪淘尽，千古风流人物。

        故垒西边，人道是、三国周郎赤壁。

        乱石穿空，惊涛拍岸，卷起千堆雪。

        江山如画，一时多少豪杰。

        \vspace{18pt}  % 分阙

        遥想公瑾当年，小乔初嫁了，雄姿英发。

        羽扇纶巾，谈笑间，樯橹灰飞烟灭。

        故国神游，多情应笑我，早生华发。

        人生如梦，一尊还酹江月。

    \end{Large}

\end{center}

\vspace{8pt}
\clearpage

\section{念奴娇·过洞庭}

\beisong{泛读}

\begin{center}
    {\zihao{-0}\bfseries 念奴娇·过洞庭}

    \vspace{18pt}

    \begin{large}\xuwen
        〔张孝祥〕
    \end{large}

\end{center}

    \vspace{18pt}  % 起首

\begin{center}
    \begin{Large}

        洞庭青草，近中秋、更无一点风色。

        玉鉴琼田三万顷，著我扁舟一叶。

        素月分辉，明河共影，表里俱澄澈。

        悠然心会，妙处难与君说。

        \vspace{18pt}  % 分阙

        应念岭海经年，孤光自照，肝胆皆冰雪。

        短发萧骚襟袖冷，稳泛沧浪空阔。

        尽吸西江，细斟北斗，万象为宾客。

        扣舷独笑，不知今夕何夕。

    \end{Large}

\end{center}

\vspace{8pt}
\clearpage

\section{贺新郎·国脉微如缕}

\beisong{精读}

\begin{center}
    {\zihao{-0}\bfseries 贺新郎·国脉微如缕}

    \vspace{18pt}

    \begin{large}\xuwen
        〔刘克庄〕
    \end{large}

\end{center}

    \vspace{18pt}  % 起首

\begin{center}
    \begin{Large}

        国脉微如缕。

        问长缨何时入手，缚将戎主？

        未必人间无好汉，谁与宽些尺度？

        试看取当年韩五。

        岂有谷城公付授，也不干曾遇骊山母。

        谈笑起，两河路。

        \vspace{18pt}  % 分阙

        少时棋柝曾联句。

        叹而今登楼揽镜，事机频误。

        闻说北风吹面急，边上冲梯屡舞。

        君莫道投鞭虚语。

        自古一贤能制难，有金汤便可无张许？

        快投笔，莫题柱。

    \end{Large}

\end{center}

\vspace{8pt}
\clearpage

\section{无家别}

\beisong{泛读}

\begin{center}
    {\zihao{-0}\bfseries 无家别}

    \vspace{18pt}

    \begin{large}\xuwen
        〔杜甫〕
    \end{large}

\end{center}

    \vspace{18pt}  % 起首

\begin{center}
    \begin{Large}

        寂寞天宝后，园庐但蒿藜。

        我里百馀家，世乱各东西。

        存者无消息，死者为尘泥。

        贱子因阵败，归来寻旧蹊。

        久行见空巷，日瘦气惨凄，

        但对狐与狸，竖毛怒我啼。

        四邻何所有，一二老寡妻。

        宿鸟恋本枝，安辞且穷栖。

        方春独荷锄，日暮还灌畦。

        县吏知我至，召令习鼓鞞。

        虽从本州役，内顾无所携。

        近行止一身，远去终转迷，

        家乡既荡尽，远近理亦齐。

        永痛长病母，五年委沟溪，

        生我不得力，终身两酸嘶，

        人生无家别，何以为蒸黎。

    \end{Large}

\end{center}

\vspace{8pt}
\clearpage

\chapter{第十二级}
\clearpage

\section{桂枝香·金陵怀古}

\beisong{背诵}

\begin{center}
    {\zihao{-0}\bfseries 桂枝香·金陵怀古}

    \vspace{18pt}

    \begin{large}\xuwen
        〔王安石〕
    \end{large}

\end{center}

    \vspace{18pt}  % 起首

\begin{center}
    \begin{Large}

        登临送目，正故国晚秋，天气初肃。

        千里澄江似练，翠峰如簇。

        征帆去棹残阳里，背西风、酒旗斜矗。

        彩舟云淡，星河鹭起，画图难足。

        \vspace{18pt}  % 分阙

        念往昔，繁华竞逐。

        叹门外楼头，悲恨相续。

        千古凭高对此，谩嗟荣辱。

        六朝旧事随流水，但寒烟衰草凝绿。

        至今商女，时时犹唱，《后庭》遗曲。

    \end{Large}

\end{center}

\vspace{8pt}
\clearpage

\section{江城子·乙卯正月二十日夜记梦}

\beisong{背诵}

\begin{center}
    {\zihao{-0}\bfseries 江城子·乙卯正月二十日夜记梦}

    \vspace{18pt}

    \begin{large}\xuwen
        〔苏轼〕
    \end{large}

\end{center}

    \vspace{18pt}  % 起首

\begin{center}
    \begin{Large}

        十年生死两茫茫。

        不思量。自难忘。

        千里孤坟，无处话凄凉。

        纵使相逢应不识，

        尘满面，鬓如霜。

        \vspace{18pt}  % 分阙

        夜来幽梦忽还乡。

        小轩窗。正梳妆。

        相顾无言，惟有泪千行。

        料得年年肠断处，

        明月夜，短松冈。

    \end{Large}

\end{center}

\vspace{8pt}
\clearpage

\section{登快阁}

\beisong{泛读}

\begin{center}
    {\zihao{-0}\bfseries 登快阁}

    \vspace{18pt}

    \begin{large}\xuwen
        〔黄庭坚〕
    \end{large}

\end{center}

    \vspace{18pt}  % 起首

\begin{center}
    \begin{Large}

        痴儿了却公家事，快阁东西倚晚晴。

        落木千山天远大，澄江一道月分明。

        朱弦已为佳人绝，青眼聊因美酒横。

        万里归船弄长笛，此心吾与白鸥盟。

    \end{Large}

\end{center}

\vspace{8pt}
\clearpage

\section{苏幕遮}

\beisong{泛读}

\begin{center}
    {\zihao{-0}\bfseries 苏幕遮}

    \vspace{18pt}

    \begin{large}\xuwen
        〔周彦邦〕
    \end{large}

\end{center}

    \vspace{18pt}  % 起首

\begin{center}
    \begin{Large}

        燎沉香，消溽暑。

        鸟雀呼晴，侵晓窥檐语。

        叶上初阳干宿雨，

        水面清圆，一一风荷举。

        \vspace{18pt}  % 分阙

        故乡遥，何日去？

        家住吴门，久作长安旅。

        五月渔郎相忆否？

        小楫轻舟，梦入芙蓉浦。

    \end{Large}

\end{center}

\vspace{8pt}
\clearpage

\section{临安春雨初霁}

\beisong{泛读}

\begin{center}
    {\zihao{-0}\bfseries 临安春雨初霁}

    \vspace{18pt}

    \begin{large}\xuwen
        〔陆游〕
    \end{large}

\end{center}

    \vspace{18pt}  % 起首

\begin{center}
    \begin{Large}

        世味年来薄似纱，谁令骑马客京华。

        小楼一夜听春雨，深巷明朝卖杏花。

        矮纸斜行闲作草，晴窗细乳戏分茶。

        素衣莫起风尘叹，犹及清明可到家。

    \end{Large}

\end{center}

\vspace{8pt}
\clearpage

\section{送灵澈}

\beisong{泛读}

\begin{center}
    {\zihao{-0}\bfseries 送灵澈}

    \vspace{18pt}

    \begin{large}\xuwen
        〔刘长卿〕
    \end{large}

\end{center}

    \vspace{18pt}

\begin{center}
    \begin{Large}

        苍苍竹林寺，杳杳钟声晚。

        荷笠带斜阳，青山独归远。

    \end{Large}

\end{center}

\vspace{8pt}
\clearpage

\section{怨情}

\beisong{泛读}

\begin{center}
    {\zihao{-0}\bfseries 怨情}

    \vspace{18pt}

    \begin{large}\xuwen
        〔李白〕
    \end{large}

\end{center}

    \vspace{18pt}

\begin{center}
    \begin{Large}

        美人卷珠帘，深坐蹙蛾眉。

        但见泪痕湿，不知心恨谁？

    \end{Large}

\end{center}

\vspace{8pt}
\clearpage

\section{贫女}

\beisong{泛读}

\begin{center}
    {\zihao{-0}\bfseries 贫女}

    \vspace{18pt}

    \begin{large}\xuwen
        〔秦韬玉〕
    \end{large}

\end{center}

    \vspace{18pt}

\begin{center}
    \begin{Large}

        蓬门未识绮罗香，拟托良媒益自伤。

        谁爱风流高格调？共怜时世俭梳妆。

        敢将十指夸针巧，不把双眉斗画长。

        苦恨年年压金线，为他人作嫁衣裳。

    \end{Large}

\end{center}

\vspace{8pt}
\clearpage

\section{菩萨蛮·书江西造口壁}

\beisong{背诵}

\begin{center}
    {\zihao{-0}\bfseries 菩萨蛮·书江西造口壁}

    \vspace{18pt}

    \begin{large}\xuwen
        〔辛弃疾〕
    \end{large}

\end{center}

    \vspace{18pt}  % 起首

\begin{center}
    \begin{Large}

        郁孤台下清江水，中间多少行人泪。

        西北望长安，可怜无数山。

        \vspace{18pt}  % 分阙

        青山遮不住，毕竟东流去。

        江晚正愁余，山深闻鹧鸪。

    \end{Large}

\end{center}

\vspace{8pt}
\clearpage

\section{长亭送别}

\beisong{泛读}

\begin{center}
    {\zihao{-0}\bfseries 长亭送别}

    \vspace{18pt}

    \begin{large}\xuwen
        〔王实甫〕
    \end{large}

\end{center}

    \vspace{18pt}  % 起首

\begin{center}
    \begin{Large}

        碧云天，黄花地，

        西风紧，北雁南飞。

        晓来谁染霜林醉？总是离人泪。

    \end{Large}

\end{center}

\vspace{8pt}
\clearpage

\section{朝天子·咏喇叭}

\beisong{泛读}

\begin{center}
    {\zihao{-0}\bfseries 朝天子·咏喇叭}

    \vspace{18pt}

    \begin{large}\xuwen
        〔王磐〕
    \end{large}

\end{center}

    \vspace{18pt}  % 起首

\begin{center}
    \begin{Large}

        喇叭，唢呐，曲儿小腔儿大。

        官船来往乱如麻，全仗你抬声价。

        \vspace{18pt}  % 分阙

        军听了军愁，民听了民怕。

        哪里去辨甚么真共假？

        \vspace{18pt}  % 分阙

        眼见的吹翻了这家，吹伤了那家，

        只吹的水尽鹅飞罢！

    \end{Large}

\end{center}

\vspace{8pt}
\clearpage

\section{一剪梅·红藕香残玉簟秋}

\beisong{泛读}

\begin{center}
    {\zihao{-0}\bfseries 一剪梅·红藕香残玉簟秋}

    \vspace{18pt}

    \begin{large}\xuwen
        〔李清照〕
    \end{large}

\end{center}

    \vspace{18pt}  % 起首

\begin{center}
    \begin{Large}

        红藕香残玉簟秋。

        轻解罗裳，独上兰舟。

        云中谁寄锦书来？

        雁字回时，月满西楼。

        \vspace{18pt}  % 分阙

        花自飘零水自流。

        一种相思，两处闲愁。

        此情无计可消除，

        才下眉头，却上心头。

    \end{Large}

\end{center}

\vspace{8pt}
\clearpage

\section{垂老别}

\beisong{泛读}

\begin{center}
    {\zihao{-0}\bfseries 垂老别}

    \vspace{18pt}

    \begin{large}\xuwen
        〔杜甫〕
    \end{large}

\end{center}

    \vspace{18pt}  % 起首

\begin{center}
    \begin{Large}

        四郊未宁静，垂老不得安。

        子孙阵亡尽，焉用身独完！

        投杖出门去，同行为辛酸。

        幸有牙齿存，所悲骨髓干。

        男儿既介胄，长揖别上官。

        老妻卧路啼，岁暮衣裳单。

        孰知是死别，且复伤其寒。

        此去必不归，还闻劝加餐。

        土门壁甚坚，杏园度亦难。

        势异邺城下，纵死时犹宽。

        人生有离合，岂择衰盛端！

        忆昔少壮日，迟回竟长叹。

        万国尽征戍，烽火被冈峦。

        积尸草木腥，流血川原丹。

        何乡为乐土？安敢尚盘桓！

        弃绝蓬室居，塌然摧肺肝。

    \end{Large}

\end{center}

\vspace{8pt}
\clearpage

\end{document}